% =====================================================================
%  LDA-T318 — Lecture 4 — Inference
%  Tuesday 10 November 2026, 16:15–17:45
%  Build:  latexmk -xelatex main.tex
% =====================================================================
\documentclass[aspectratio=169,10pt]{beamer}

% Path has no parent-directory step: Overleaf resolves \input from the project root and
% will not read above it. Locally, the Makefile puts common/ on TEXINPUTS.
% =====================================================================
%  LDA-T318 — Large Language Models and Generative AI for NLP
%  Shared Beamer preamble for all lecture decks.
%  Aarne Talman, University of Helsinki, 2026.
%
%  Engine:  XeLaTeX or LuaLaTeX preferred (gives the Fira font set).
%           pdfLaTeX also works; Metropolis falls back to Latin Modern.
%  Build:   latexmk -xelatex main.tex      (from the deck's own directory)
%
%  Every deck does:
%      \documentclass[aspectratio=169,10pt]{beamer}
%      \input{../common/preamble.tex}
% =====================================================================

% ---------------------------------------------------------------- theme
\usetheme{metropolis}
\metroset{
  titleformat=smallcaps,
  numbering=fraction,
  progressbar=frametitle,
  block=fill,
  sectionpage=progressbar,
  subsectionpage=none,
}

% ---------------------------------------------------------------- packages
% inputenc/fontenc are required under pdfLaTeX and actively harmful under
% XeLaTeX/LuaLaTeX, where Metropolis loads fontspec to get Fira. Guard them.
\usepackage{iftex}
\ifPDFTeX
  \usepackage[utf8]{inputenc}
  \usepackage[T1]{fontenc}
  \usepackage[sfdefault]{FiraSans}
  \usepackage{FiraMono}
\else
  % Metropolis only gets Fira if fontspec can actually find it. TeX Live ships
  % the OTFs but does NOT register them with fontconfig, so a plain
  % \setsansfont{Fira Sans} silently falls back to Latin Modern and the deck
  % stops looking like Metropolis. Load them by filename instead.
  \usepackage{fontspec}
  \setsansfont{FiraSans}[
    Extension      = .otf,
    UprightFont    = *-Book,
    ItalicFont     = *-BookItalic,
    BoldFont       = *-Medium,
    BoldItalicFont = *-MediumItalic]
  \setmonofont{FiraMono}[
    Extension   = .otf,
    UprightFont = *-Regular,
    BoldFont    = *-Medium,
    Scale       = 0.92]
  \renewcommand{\familydefault}{\sfdefault}
\fi
\usepackage[english]{babel}

\usepackage{amsmath,amssymb}
\usepackage{booktabs}
\usepackage{array}
\usepackage{multicol}
\usepackage{xcolor}
\usepackage{graphicx}
\usepackage{tikz}
\usepackage{pgfplots}
\usepackage{listings}
\usepackage{appendixnumberbeamer}

\pgfplotsset{compat=1.17}
\usetikzlibrary{positioning,arrows.meta,fit,backgrounds,calc,matrix,shapes.geometric,decorations.pathreplacing}

% =====================================================================
%  University of Helsinki palette
%  Primary brand colour is UH Hray/black with the "HY blue" accent.
%  Secondary accents chosen for projector contrast, not for print.
% =====================================================================
\definecolor{UHBlue}{HTML}{0C5A9E}   % primary accent
\definecolor{UHDark}{HTML}{15202B}   % body text / dark backgrounds
\definecolor{UHGrey}{HTML}{5A6672}   % secondary text
\definecolor{UHLight}{HTML}{EEF2F6}  % block backgrounds
\definecolor{UHOrange}{HTML}{D1601A} % alerts, gotchas
\definecolor{UHGreen}{HTML}{1F7A5C}  % examples, "this works"
\definecolor{UHRed}{HTML}{A62B2B}    % failure modes
\definecolor{UHSand}{HTML}{C8B88A}   % neutral third accent

\setbeamercolor{normal text}{fg=UHDark, bg=white}
\setbeamercolor{alerted text}{fg=UHOrange}
\setbeamercolor{example text}{fg=UHGreen}
\setbeamercolor{frametitle}{fg=white, bg=UHBlue}
\setbeamercolor{title separator}{fg=UHBlue}
\setbeamercolor{progress bar}{fg=UHBlue}
\setbeamercolor{progress bar in head/foot}{fg=UHBlue}
\setbeamercolor{progress bar in section page}{fg=UHBlue}
\setbeamercolor{structure}{fg=UHBlue}
\setbeamercolor{palette primary}{fg=white, bg=UHBlue}   % section pages, standout frames
\setbeamercolor{block title}{fg=white, bg=UHBlue}
\setbeamercolor{block body}{fg=UHDark, bg=UHLight}
\setbeamercolor{block title alerted}{fg=white, bg=UHOrange}
\setbeamercolor{block body alerted}{fg=UHDark, bg=UHLight}
\setbeamercolor{block title example}{fg=white, bg=UHGreen}
\setbeamercolor{block body example}{fg=UHDark, bg=UHLight}
\setbeamercolor{footnote}{fg=UHGrey}

\setbeamerfont{frametitle}{size=\large, series=\bfseries}
\setbeamerfont{footnote}{size=\tiny}
\setbeamerfont{caption}{size=\scriptsize}

% Frame title bar.
%
% Do NOT reimplement the frametitle template here. Metropolis's own template
% already spans \paperwidth (so `bg=UHBlue' above gives the full-width bar),
% already struts two-line titles without collision, and already draws the
% frametitle progress bar. Reimplementing it means touching internals whose
% names contain `@', which is NOT a letter in a .tex file -- `\@empty' parses
% as `\@' followed by the word "empty" and the build breaks.
%
% To change the bar's height, adjust Metropolis's own padding length instead,
% inside \makeatletter so that `@' is a letter:
%
%   \makeatletter
%   \setlength{\metropolis@frametitle@padding}{2.4ex}  % default 2.2ex
%   \makeatother
%
\makeatletter
\setlength{\metropolis@frametitle@padding}{1.2ex}
\makeatother

% No deck uses \framesubtitle; if one ever does, Metropolis ignores it and the
% subtitle should go in the frame body instead.

% =====================================================================
%  Course metadata — decks override \lecturename and \lecturenumber
% =====================================================================
\newcommand{\coursecode}{LDA-T318}
\newcommand{\coursename}{LLMs and Generative AI for NLP}
\newcommand{\courseteacher}{Aarne Talman}
\newcommand{\courseinstitution}{University of Helsinki}
\providecommand{\lecturename}{}
\providecommand{\lecturedate}{}

% Footline: course code on the left, lecture name centred, n/N on the right.
\setbeamertemplate{frame footer}{\coursecode\quad\textbar\quad\lecturename}

% =====================================================================
%  Semantic macros — use these instead of ad-hoc formatting, so that a
%  change of house style is a single edit here.
% =====================================================================

% --- Sourcing. Every quantitative claim on a slide should carry one. ---
\newcommand{\source}[1]{%
  \vfill{\scriptsize\color{UHGrey}\raggedright Source: #1\par}%
}
\newcommand{\sources}[1]{%
  \vfill{\scriptsize\color{UHGrey}\raggedright Sources: #1\par}%
}
% Inline citation chip, e.g. \cit{Raschka 2026a}
\newcommand{\cit}[1]{{\scriptsize\color{UHGrey}[#1]}}

% --- Counter-intuitive point. Three per lecture; these are exam material. ---
\newenvironment{gotcha}[1][Gotcha]{%
  \begin{alertblock}{#1}%
}{%
  \end{alertblock}%
}

% --- Speculative / unconfirmed claim. Research literacy matters here. ---
\newcommand{\rumour}{{\color{UHOrange}\textbf{[unconfirmed]}}\ }

% --- Key takeaway, for the closing slide of each section. ---
\newenvironment{takeaway}{%
  \begin{block}{Take away}%
}{%
  \end{block}%
}

% --- Live demo marker. ---
\newcommand{\demo}[1]{{\color{UHGreen}\textbf{Demo:}} #1}

% --- Emphasis consistent with the palette. ---
\newcommand{\key}[1]{{\color{UHBlue}\textbf{#1}}}
\newcommand{\warn}[1]{{\color{UHOrange}\textbf{#1}}}
\newcommand{\good}[1]{{\color{UHGreen}\textbf{#1}}}
\newcommand{\bad}[1]{{\color{UHRed}\textbf{#1}}}

% --- Units and quantities used constantly across the decks. ---
\newcommand{\unit}[1]{\,\mathrm{#1}}
\newcommand{\flops}{\mathrm{FLOPs}}
% NB: no \toks macro here --- \toks is a TeX primitive (a token register), so
% \newcommand{\toks} is a hard "already defined" error. No deck used it.
\newcommand{\KV}{\mathrm{KV}}
\newcommand{\softmax}{\operatorname{softmax}}

% =====================================================================
%  Code listings — used for config snippets, Slurm scripts, chat templates.
%  Frames containing these MUST be declared [fragile].
% =====================================================================
\lstdefinestyle{uh}{
  basicstyle=\ttfamily\scriptsize\color{UHDark},
  keywordstyle=\color{UHBlue}\bfseries,
  commentstyle=\color{UHGrey}\itshape,
  stringstyle=\color{UHGreen},
  numberstyle=\tiny\color{UHGrey},
  backgroundcolor=\color{UHLight},
  frame=none,
  xleftmargin=0.6em,
  framexleftmargin=0.6em,
  breaklines=true,
  breakatwhitespace=true,
  showstringspaces=false,
  columns=fullflexible,
  keepspaces=true,
  upquote=true,
  literate={-}{{-}}1 {_}{{\_}}1,
}
\lstset{style=uh}

\lstdefinelanguage{slurm}{
  morekeywords={SBATCH,module,srun,source,export,set,echo},
  morecomment=[l]{\#},
  sensitive=true,
}
\lstdefinelanguage{jinja}{
  morekeywords={for,endfor,if,endif,set,in},
  morecomment=[s]{\{\#}{\#\}},
  sensitive=true,
}

% =====================================================================
%  TikZ styles shared across decks (block diagrams, token grids).
% =====================================================================
\tikzset{
  mod/.style     = {draw=UHBlue, fill=UHLight, rounded corners=2pt,
                    minimum height=7mm, minimum width=22mm,
                    font=\scriptsize, align=center},
  modhi/.style   = {mod, draw=UHOrange, fill=UHOrange!12},
  modmute/.style = {mod, draw=UHGrey, fill=white, text=UHGrey},
  flow/.style    = {-{Latex[length=2mm]}, draw=UHGrey, thick},
  flowhi/.style  = {flow, draw=UHOrange},
  resid/.style   = {-{Latex[length=2mm]}, draw=UHBlue, thick, dashed},
  tokcell/.style = {draw=UHGrey!50, minimum size=4.2mm, font=\tiny, inner sep=0pt},
  tokon/.style   = {tokcell, fill=UHBlue!35},
  tokoff/.style  = {tokcell, fill=UHLight},
  tokmask/.style = {tokcell, fill=UHGrey!15, text=UHGrey!60},
  lbl/.style     = {font=\scriptsize\color{UHGrey}},
}

% =====================================================================
%  Presenter notes.
%
%  Notes are written on every slide. By default they are NOT shown, so
%  `latexmk -xelatex main.tex` produces a clean deck for projection.
%
%  To present with notes on a second screen (pdfpc, Skim, Impressive):
%      uncomment the line below, rebuild to main-notes.pdf
%  To print a notes handout:
%      \setbeameroption{show only notes}
% =====================================================================
% \setbeameroption{show notes on second screen=right}

\setbeamertemplate{note page}{%
  {\scriptsize\color{UHGrey}\coursecode\ \textbar\ \lecturename\par}%
  \vskip0.3em\hrule\vskip0.6em%
  {\small\insertnote}%
}

% ---------------------------------------------------------------- misc
\setbeamertemplate{caption}[numbered]
\setbeamertemplate{itemize items}[circle]
\setlength{\parskip}{0.1em}

% ---------------------------------------------------------------- density
% These decks are text-dense and a 16:9 frame is only 90mm tall, so the
% default vertical rhythm costs several lines per slide. Tightened globally
% here rather than slide by slide.
\setbeamersize{text margin left=0.7cm, text margin right=0.7cm}
\makeatletter
\def\@listi{\leftmargin\leftmargini
  \topsep 1.5\p@ \@plus1\p@ \@minus1\p@
  \parsep 0\p@
  \itemsep 1.5\p@ \@plus1\p@ \@minus1\p@}
\let\@listI\@listi
\makeatother

% The title page is a stack of short paragraphs, so a non-zero \parskip
% pushes it past the bottom of the frame. Neutralise it there only.
\addtobeamertemplate{title page}{\setlength{\parskip}{0pt}}{}

% Tighter lists — the decks are dense.
\setbeamertemplate{itemize/enumerate body begin}{\small}
\setbeamertemplate{itemize/enumerate subbody begin}{\footnotesize}

% Table helpers
\newcolumntype{L}[1]{>{\raggedright\arraybackslash}p{#1}}
\newcolumntype{C}[1]{>{\centering\arraybackslash}p{#1}}
\newcolumntype{R}[1]{>{\raggedleft\arraybackslash}p{#1}}

% Allow slightly overfull boxes without a warning storm
\hfuzz=2pt

% =====================================================================
%  Shared course-map frame.
%  Every lecture opens with this, with the current week highlighted:
%      \coursemapframe{1}   ... \coursemapframe{4}
%  Pass 0 to highlight nothing.
% =====================================================================
% \cmrow{current week}{this row's week}{dates}{Tuesday}{Thursday}{teacher}
% Highlights the row in UH blue bold when the two week numbers match.
% Deliberately avoids \rowcolor so that colortbl is not needed.
% The & and \\ must NOT sit inside the \ifnum branches: TeX scans alignment
% cells while skipping the false branch, so an & there breaks the conditional
% ("Incomplete \ifnum"). Keep the branches down to a macro name and let the
% separators live in the macro bodies, which expand after the test resolves.
\newcommand{\cmhi}[5]{%
  \textcolor{UHBlue}{\textbf{#1}} & \textcolor{UHBlue}{\textbf{#2}} &
  \textcolor{UHBlue}{\textbf{#3}} & \textcolor{UHBlue}{\textbf{#4}} &
  \textcolor{UHBlue}{\textbf{#5}}\\%
}
\newcommand{\cmplain}[5]{#1 & #2 & #3 & #4 & #5\\}
\newcommand{\cmrow}[6]{%
  \ifnum#1=#2\relax\expandafter\cmhi\else\expandafter\cmplain\fi
  {#2}{#3}{#4}{#5}{#6}%
}
\newcommand{\coursemapframe}[1]{%
\begin{frame}{Course overview}
\centering
\small
\setlength{\parskip}{0pt}
\renewcommand{\arraystretch}{1.08}
\begin{tabular}{@{}c L{1.7cm} L{3.1cm} L{3.7cm} L{1.4cm}@{}}
\toprule
\textbf{Wk} & \textbf{Dates} & \textbf{Tuesday} & \textbf{Thursday} & \textbf{Teacher}\\
\midrule
\cmrow{#1}{1}{27 / 29 Oct}{Language models}{Pretraining}{Aarne}
\cmrow{#1}{2}{3 / 5 Nov}{Post-training}{Lab 2: tune vs prompt}{Aarne}
\cmrow{#1}{3}{10 / 12 Nov}{Inference}{Lab 3: serve}{Aarne}
\cmrow{#1}{4}{17 / 19 Nov}{Evaluation}{Lab 4: evaluate}{Jussi}
\cmrow{#1}{5}{24 / 26 Nov}{RAG}{Lab 5: RAG}{Dmitry}
\cmrow{#1}{6}{1 / 3 Dec}{Agents}{Lab 6: agent}{Dmitry}
\cmrow{#1}{7}{8 / 10 Dec}{Course recap}{Special topics (TBC)}{All}
\bottomrule
\end{tabular}

\vspace{1.2em}
{\footnotesize Reusing a task and model across labs is recommended, not required.\\
There are no intermediate submissions; the final report and code repository
determine the grade.}
\note{Fifteen seconds, not more. The point of this slide is orientation, not
content: show them the shape of the course and where today sits inside it.
The final week is deliberately left flexible for a recap and special topics.}
\end{frame}%
}


\renewcommand{\lecturename}{Lecture 4 — Inference}
\renewcommand{\lecturedate}{10 November 2026}

\title{Inference}
\subtitle{Lecture 4: what it costs to run a language model}
\author{Aarne Talman}
\institute{\coursecode\ \textbar\ \coursename\ \textbar\ \courseinstitution}
\date{10 November 2026}

\begin{document}

% =====================================================================
\begin{frame}[plain,shrink=8]
  \titlepage
  \note{Set the frame in the first minute. Weeks 1 and 2 were about where the
  weights come from and how to shape their behaviour. Today's question is what
  the finished model costs to run, every hour, for as long as it is deployed.

  The one sentence for the whole lecture: you pay the training cost once and
  the inference cost every hour after that. The technical reason it all hangs
  together is that decoding is memory-bandwidth bound, so almost every serving
  trick is a way to move fewer bytes per generated token, or to spread the
  bytes you do move across more users.

  Thursday they serve their own model and measure it. Several numbers on these
  slides are estimates they will check against a stopwatch in 48 hours. Tell
  them that, because it changes how they listen.}
\end{frame}

\coursemapframe{3}

\begin{frame}{Lecture overview}
  \begin{columns}[T]
    \begin{column}{0.5\textwidth}
      \begin{enumerate}
        \item Decoding \hfill \textit{10 min}
        \item Prefill vs decode \hfill \textit{10 min}
        \item The KV cache \hfill \textit{15 min}
        \item Continuous batching \hfill \textit{10 min}
      \end{enumerate}
    \end{column}
    \begin{column}{0.5\textwidth}
      \begin{enumerate}
        \setcounter{enumi}{4}
        \item Quantisation for inference \hfill \textit{10 min}
        \item Speculative and constrained decoding \hfill \textit{10 min}
        \item Serving systems and cost \hfill \textit{15 min}
        \item Lab 3: measuring a server fairly \hfill \textit{10 min}
      \end{enumerate}
    \end{column}
  \end{columns}
  \vfill
  \begin{takeaway}
    By the end of today you should be able to explain where inference time and
    memory go, estimate a model's serving capacity, and design a fair
    speed--quality comparison.
  \end{takeaway}
  \note{Read the timings out loud as in every lecture. Eight sections is a lot,
  so the signposting matters more than usual today.

  The takeaway box is the promise. Come back to it at 17:43. Note that it is a
  promise about \emph{paper} estimates. The measurement happens on Thursday,
  and the gap between the two is the most educational part of the lab.}
\end{frame}

% =====================================================================
\section{Decoding}
% =====================================================================

\begin{frame}{From logits to a token}
  Every generated token goes through the same short pipeline. All of it
  happens \emph{after} the forward pass, on a vector of
  $|V| \approx 150{,}000$ numbers.
  \vspace{0.6em}
  \centering
  \begin{tikzpicture}[node distance=5mm]
    \node[modmute] (lg) {logits\\{\tiny $z \in \mathbb{R}^{|V|}$}};
    \node[mod, right=of lg] (pen) {penalties\\{\tiny repetition, presence}};
    \node[modhi, right=of pen] (tmp) {temperature\\{\tiny $z/T$}};
    \node[modhi, right=of tmp] (trn) {truncation\\{\tiny top-$k$, top-$p$, min-$p$}};
    \node[mod, right=of trn] (sm) {softmax\\{\tiny $\to$ probabilities}};
    \node[modmute, below=8mm of sm] (dr) {draw one token};
    \draw[flow] (lg) -- (pen);
    \draw[flow] (pen) -- (tmp);
    \draw[flow] (tmp) -- (trn);
    \draw[flow] (trn) -- (sm);
    \draw[flow] (sm) -- (dr);
  \end{tikzpicture}

  \vspace{0.6em}
  \raggedright
  Two things worth noticing now:
  \begin{itemize}
    \item it is \key{cheap}: a few operations over one vector, next to
      billions of FLOPs in the forward pass;
    \item it is \key{not part of the model}. The same weights behave very
      differently under different settings, and those settings live in
      \texttt{generation\_config.json}, in your client call, or in your
      server's defaults.
  \end{itemize}
  \note{Clock check: should be about 16:18 when this goes up.

  The ordering matters and people get it wrong: penalties and temperature act
  on \emph{logits}; truncation acts on the sorted distribution. Divide by $T$
  before truncating, and top-$p$ sees a different distribution than if you
  truncate first. Implementations differ, which is one small reason two
  libraries disagree at the same nominal settings.

  The second bullet is the one to land. Three of the five boxes sit outside the
  \texttt{.safetensors} file, which is the same point as the ``what a model
  actually is'' slide in week 1.}
\end{frame}

\begin{frame}{Temperature and truncation change one distribution}
  \centering
  \begin{tikzpicture}
    \pgfplotsset{every axis/.append style={
      scale only axis, width=2.0cm, height=2.1cm,
      ybar, bar width=2.6pt,
      xmin=0.4, xmax=5.6, ymin=0, ymax=1.05,
      ytick={0,0.5,1}, xtick={1,2,3,4,5},
      xticklabels={a,b,c,d,e},
      tick label style={font=\tiny},
      title style={font=\tiny, yshift=-0.5mm},
      axis x line*=bottom, axis y line*=left,
    }}
    \begin{axis}[name=g0, title={raw $p$}]
      \addplot[fill=UHGrey!30, draw=UHGrey] coordinates
        {(1,0.40)(2,0.25)(3,0.18)(4,0.10)(5,0.07)};
    \end{axis}
    \begin{axis}[name=g1, at={($(g0.east)+(7mm,0)$)}, anchor=west,
                 title={greedy}]
      \addplot[fill=UHBlue!45, draw=UHBlue] coordinates
        {(1,1.00)(2,0)(3,0)(4,0)(5,0)};
    \end{axis}
    \begin{axis}[name=g2, at={($(g1.east)+(7mm,0)$)}, anchor=west,
                 title={temp 1.5}]
      \addplot[fill=UHBlue!45, draw=UHBlue] coordinates
        {(1,0.33)(2,0.24)(3,0.19)(4,0.13)(5,0.10)};
    \end{axis}
    \begin{axis}[name=g3, at={($(g2.east)+(7mm,0)$)}, anchor=west,
                 title={top-$k$, 2}]
      \addplot[fill=UHBlue!45, draw=UHBlue] coordinates
        {(1,0.62)(2,0.38)(3,0)(4,0)(5,0)};
    \end{axis}
    \begin{axis}[name=g4, at={($(g3.east)+(7mm,0)$)}, anchor=west,
                 title={top-$p$, 0.8}]
      \addplot[fill=UHBlue!45, draw=UHBlue] coordinates
        {(1,0.48)(2,0.30)(3,0.22)(4,0)(5,0)};
    \end{axis}
  \end{tikzpicture}

  \vspace{0.3em}
  \raggedright\footnotesize
  \begin{tabular}{@{}L{2.4cm} L{9.0cm}@{}}
    \textbf{greedy} & take the most likely token.\\
    \textbf{temperature} & values below 1 sharpen the distribution; values above 1 flatten it.\\
    \textbf{top-$k$} & keep a fixed number of likely tokens, then renormalise.\\
    \textbf{top-$p$} & keep the smallest set whose probabilities add up to $p$; the set adapts to confidence.\\
  \end{tabular}
  \note{Walk left to right across the little plots; they are all the same five
  candidate tokens. The two that matter conceptually are top-$p$ and min-$p$,
  because both are \emph{adaptive}: the number of surviving tokens depends on
  how confident the model is, which is what you want.

  Symbolic forms, if anyone wants them on the board: greedy is $\arg\max_i z_i$,
  the limit of $T\to 0$; temperature is the map $z \mapsto z/T$; top-$p$ keeps
  the smallest prefix of the sorted distribution with $\sum p \ge p$; min-$p$
  keeps $\{i : p_i \ge \alpha\, p_{\max}\}$ with $\alpha$ typically 0.05--0.2.

  If anyone asks about min-$p$, typical-$p$, eta or epsilon sampling, or DRY:
  they are all real, but outside today's overview.}
\end{frame}

\begin{frame}{Which sampler for which task}
  \centering
  \small
  \begin{tabular}{@{}L{3.3cm} L{3.4cm} L{4.6cm}@{}}
    \toprule
    \textbf{Task} & \textbf{Sensible default} & \textbf{Why}\\
    \midrule
    Classification, extraction, exact-match QA &
      temperature 0 (greedy) &
      there is one right answer, so variance only costs you\\[0.3em]
    Code, maths, anything verified &
      temperature 0 for one shot; about 0.8 if you sample several and vote &
      diversity only pays when you have a verifier to pick the winner\\[0.3em]
    Open generation, dialogue, summarisation &
      temperature 0.7, with top-$p$ 0.9 or min-$p$ 0.05 &
      greedy output is repetitive and flat, and a human is judging\\[0.3em]
    \bottomrule
  \end{tabular}

  \vspace{0.5em}
  \raggedright
  \key{Sample widely only when a human or an automatic checker can reject poor
  outputs.} Otherwise sampling adds variance without a way to select the
  better result.
  \source{Raschka, \textit{Build Your Own Reasoning Model}, Oct 2026,
    on best-of-$n$ and self-consistency.}
  \note{The organising sentence is the last one on the slide, so make them
  write it down. Which sampler is best depends on the \emph{pipeline} the model
  sits in.

  For their projects: three of the four suggested tasks are exact-match or F1
  against a label, so they should be running greedy and holding it fixed.
  Someone every year reports a two-point improvement that turns out to be
  sampling noise at $T=1$.}
\end{frame}

\begin{frame}{Decoding parameters are part of your experimental setup}
  \begin{columns}[T]
    \begin{column}{0.54\textwidth}
      Three places a sampling parameter can come from, in increasing order of
      surprise:
      \begin{enumerate}
        \item your API call
        \item the server's defaults
        \item the model's own \texttt{generation\_config.json}, shipped in the
          repository and silently applied by \texttt{transformers}
      \end{enumerate}
      \vspace{0.5em}
      They disagree. \texttt{transformers} honours the model's config; a server
      may or may not; OpenAI-compatible clients have their own defaults for
      \texttt{temperature} and \texttt{top\_p}.
    \end{column}
    \begin{column}{0.46\textwidth}
      \begin{block}{The rule for weeks 4--6}
        \small
        Set decoding parameters \textbf{explicitly}, in one place, and record
        them next to every number you report.
        \vspace{0.4em}
        A speed-versus-quality table is meaningless if the two rows were
        generated at different temperatures, and that is the easiest mistake
        to make by accident in Thursday's lab.
      \end{block}
      \vspace{0.3em}
      \small
      Also fix \texttt{max\_tokens} and the random seed, and report the
      \emph{distribution} of your measurements, not only the mean.
    \end{column}
  \end{columns}
  \note{Short, unglamorous slide that saves people a lot of time. Jussi will
  return to the same point in week 4 from the evaluation side. This is the
  serving side of it.

  Concrete: a model whose \texttt{generation\_config.json} sets
  \texttt{temperature=0.7} will not be greedy just because you did not ask for
  sampling. People lose an afternoon to this.

  Foreshadow section 8 lightly: even with everything pinned, including the
  seed, two runs can differ. We will find out why at the end.}
\end{frame}

\begin{frame}{Reasoning on or off: what thinking tokens cost}
  A reasoning model pays for its accuracy in \key{output tokens}, which are the
  expensive kind, because each one costs a full forward pass.
  \vspace{0.5em}
  \begin{columns}[T]
    \begin{column}{0.52\textwidth}
      \small
      \begin{itemize}
        \item Effort level correlates with token usage, and token usage
          correlates with accuracy \cit{Raschka, Jul 2026}
        \item Kimi's \textit{Toggle} training cuts generated tokens by
          \textbf{25--30\%} with little benchmark change, and the saving
          transfers from maths and code RL to GPQA and MMLU-Pro
        \item GPT-5.6 \textit{Luna} uses \textbf{80\% more tokens} than
          \textit{Sol} at comparable scores: a smaller model buys the same
          result with more thinking
      \end{itemize}
    \end{column}
    \begin{column}{0.48\textwidth}
      \begin{block}{Worked, for a serving budget}
        \small
        Answer without thinking: 150 output tokens.\\
        Answer with thinking: 1{,}800 output tokens.\\[0.3em]
        At a per-user TPOT of 5\,ms that is
        \textbf{0.75\,s} versus \textbf{9\,s} of GPU-seconds for the
        \emph{same question}, a $12\times$ difference in both latency and cost.
      \end{block}
      \vspace{0.2em}
      \small
      And the KV cache grows with every one of those tokens, so a reasoning
      request also occupies more memory for longer.
    \end{column}
  \end{columns}
  \sources{Raschka, \textit{Controlling Reasoning Effort in LLMs}, Jul 2026;
    Kimi K2.5 / Toggle report; Artificial Analysis Intelligence Index v4.3,
    as reported in Raschka, Jul 2026.}
  \note{Say the $12\times$ slowly. Turning reasoning on is an
  order-of-magnitude decision about your serving bill, not a small quality
  knob.

  Second-order point that will matter in section 3: thinking tokens are KV
  cache. A reasoning request holds a GPU slot for twelve times as long, so your
  concurrency drops by the same factor even though nothing about the model
  changed.

  Date everything here. These are late-2026 figures about specific models and
  they will be wrong by next autumn.}
\end{frame}

% Model-specific controls change quickly and are retained as optional reference.
\iffalse
\begin{frame}[shrink=4]{Asking for less thinking: the interfaces, late 2026}
  \centering
  \footnotesize
  \begin{tabular}{@{}L{2.5cm} L{4.6cm} L{3.6cm}@{}}
    \toprule
    \textbf{Model} & \textbf{How you ask} & \textbf{Values}\\
    \midrule
    gpt-oss, GPT-5.6 & system prompt \texttt{"Reasoning effort: ..."} & low / medium / high\\
    Qwen3 & tokenizer \texttt{enable\_thinking}; \texttt{/think}, \texttt{/no\_think} flags & on / off, plus a hard token budget\\
    GLM-5 & chat template, turn-level \texttt{thinking: disabled} & on / off\\
    Nemotron 3 Ultra & chat template, plus an independent token budget & off / medium / regular\\
    DeepSeek V4 & three trained modes & Non-think / Think High / Think Max\\
    Kimi K3 & \texttt{reasoning\_effort} API parameter & low / high / max (default)\\
    Inkling & system message \texttt{"Thinking effort level: 0.8"} & continuous in $[0,1]$\\
    \bottomrule
  \end{tabular}

  \vspace{0.7em}
  \raggedright
  \small
  Recall from week 2: these instructions \textbf{select a behaviour that
  post-training created}. Pasting \texttt{"Reasoning effort: high"} into an
  arbitrary model does nothing at all.
  \source{Raschka, \textit{Controlling Reasoning Effort in LLMs}, Jul 2026.}
  \note{Do not read the table. Point at three rows: the system-prompt route
  (gpt-oss), the chat-template route (Qwen3, GLM-5, Nemotron) and the
  API-parameter route (Kimi K3). Then point at Inkling, which is the odd one
  out because its control is continuous rather than ordinal.

  The callback to Lecture 3 is worth repeating: these are interfaces to a
  behaviour that post-training put there.

  Serving consequence, which is why it is in \emph{this} lecture: effort is a
  per-request parameter, so your server has to cope with requests whose output
  lengths differ by an order of magnitude, which is exactly the situation that
  makes static batching fall apart in section 4.}
\end{frame}

\begin{frame}{The hook: ``which model'' and ``how much thinking'' are one question}
  \centering
  \begin{tikzpicture}
    \begin{semilogxaxis}[
      width=7.6cm, height=4.6cm,
      xlabel={cost per task (relative, log scale)},
      ylabel={benchmark score},
      xmin=0.2, xmax=100, ymin=45, ymax=88,
      xtick={0.3,1,3,10,30,100},
      xticklabels={0.3,1,3,10,30,100},
      tick label style={font=\tiny}, label style={font=\scriptsize},
      grid=both, grid style={UHGrey!20},
      legend style={font=\tiny, at={(0.97,0.03)}, anchor=south east,
                    draw=none, fill=white},
    ]
      \addplot[UHBlue, thick, mark=*, mark size=1.4pt] coordinates
        {(0.3,52) (0.8,61) (2.0,68) (5.0,72) (12,74)};
      \addlegendentry{smaller model, effort low \textrightarrow{} max}
      \addplot[UHOrange, thick, mark=square*, mark size=1.4pt] coordinates
        {(1.5,63) (4.0,71) (10,77) (25,80) (60,81)};
      \addlegendentry{larger model, effort low \textrightarrow{} max}
      \draw[UHGrey, dashed] (axis cs:4.0,45) -- (axis cs:4.0,88);
      \node[lbl, anchor=south west] at (axis cs:4.3,46) {same cost, same score};
    \end{semilogxaxis}
  \end{tikzpicture}

  \vspace{0.2em}
  \raggedright\small
  The curves \key{overlap}. A smaller model at high reasoning effort reaches
  roughly the score of a larger model at low effort, for roughly the same
  money. Both curves are also \warn{flat at the right-hand end}, so past some
  point you pay more and learn nothing.
  \source{Shape after the Artificial Analysis cost-versus-index curves
    discussed in Raschka, \textit{Controlling Reasoning Effort}, Jul 2026.
    \textbf{Schematic: the shape is the claim, not the values.}}
  \note{This is the slide that reframes the whole lecture, so give it a full
  ninety seconds.

  The naive mental model is two separate decisions: pick a model, then pick
  settings. The curves say it is \emph{one budget decision} along a single
  cost axis, and that the two levers are partly interchangeable.

  Two caveats to state explicitly, because this is a research-literacy course.
  First, the x-axis is API \emph{price} rather than compute, so it mixes
  provider margins into a technical plot. Second, the curve shapes are
  benchmark-specific; this overlap is reported for coding-agent indices and
  should not be assumed everywhere.

  Close the section: everything after this is about moving your own curve left,
  i.e. getting the same score for less money, on hardware you control.}
\end{frame}
\fi

% =====================================================================
\section{Prefill vs decode}
% =====================================================================

\begin{frame}{One request, two different computations}
  \centering
  \begin{tikzpicture}[x=1mm, y=1mm]
    % prefill block
    \fill[UHBlue!35, draw=UHBlue] (0,0) rectangle (26,7);
    \node[font=\scriptsize] at (13,3.5) {prefill};
    % decode steps
    \foreach \i in {0,...,9} {
      \fill[UHOrange!25, draw=UHOrange] ({28+\i*7},0) rectangle ({33+\i*7},7);
    }
    \node[font=\scriptsize] at (62,3.5) {};
    \node[lbl, anchor=west] at (99,3.5) {decode, one token per step};
    % arrows
    \draw[flow] (0,-4) -- (26,-4);
    \node[lbl, anchor=north] at (13,-5) {TTFT};
    \draw[flow] (28,12) -- (35,12);
    \node[lbl, anchor=south] at (31.5,12.5) {TPOT};
    \node[lbl, anchor=east] at (-2,3.5) {time $\rightarrow$};
  \end{tikzpicture}
  \vspace{0.3em}

  \centering
  \small
  \begin{tabular}{@{}L{2.2cm} L{4.4cm} L{4.6cm}@{}}
    \toprule
    & \textbf{Prefill} & \textbf{Decode}\\
    \midrule
    What & the whole prompt at once & one new token per step\\
    Shape & matrix $\times$ matrix & matrix $\times$ \emph{vector}\\
    Reads & every weight once, for 2{,}000 tokens & every weight once, for 1 token\\
    Limited by & \key{how fast the chip multiplies} & \key{how fast memory feeds it}\\
    Felt as & \key{TTFT} --- wait for the first token & \key{TPOT} --- gap between tokens\\
    \bottomrule
  \end{tabular}
  \note{Clock check: should be around 16:28.

  Draw the picture on the board as well. This asymmetry matters more than
  anything else in serving, and it follows directly from the causal
  mask in week 1: training and prefill can process the whole triangle in
  parallel, decode cannot, because token $t+1$ does not exist until token $t$
  has been produced.

  The line to repeat: in decode you read the entire model from HBM to produce
  \emph{one} token. That is an absurd ratio, and every technique in the rest of
  the lecture is a reaction to it.}
\end{frame}

\begin{frame}{Fetching a number costs far more than multiplying it}
  A GPU does two things, and they are nothing like equally expensive:
  \begin{itemize}
    \item \key{fetching} a weight out of the card's memory into the arithmetic
      units --- \textbf{slow};
    \item \key{multiplying} with it once it has arrived --- roughly a hundred
      times cheaper.
  \end{itemize}
  So the number that decides your speed is: \emph{how much arithmetic do you
  get out of each byte you went to the trouble of fetching?} That ratio is
  called \key{arithmetic intensity}, measured in FLOPs per byte read.
  \vspace{0.3em}
  \begin{columns}[T]
    \begin{column}{0.5\textwidth}
      \begin{block}{Decode: one token at a time}
        \small
        Read all 8\,GB of weights and use each weight in one multiply-add per
        token. Batch 32 users and each fetched weight serves 32 tokens.
        Intensity $\approx$ the number of tokens in the step.
      \end{block}
    \end{column}
    \begin{column}{0.5\textwidth}
      \begin{block}{Prefill: the whole prompt at once}
        \small
        A 2{,}000-token prompt goes through in one pass. The same 8\,GB is read
        once and reused two thousand times over. Intensity $\approx 2000$.
      \end{block}
    \end{column}
  \end{columns}
  \vspace{0.4em}
  \centering
  \key{For ordinary decode, batching is the main way to reuse fetched
  weights.}
  \note{The derivation, which is worth doing on the board because it explains
  everything that follows. Write $I = \text{FLOPs} / \text{bytes moved}$. For
  the linear layers of a dense model with $N$ parameters in bf16, running $B$
  tokens through one forward pass: $\text{FLOPs} \approx 2NB$ (each weight does
  one multiply-add, i.e.\ 2 FLOPs, per token), $\text{bytes read} \approx 2N$
  (two bytes per weight, read once). Hence $I \approx B$ FLOP per byte, so the
  arithmetic intensity is just the batch size. Decode with one sequence is
  $B=1$; batching 32 users gives $I \approx 32$; a 2{,}000-token prefill is
  $B = 2000$.

  Note the deliberate simplification: this ignores attention over the KV cache,
  which at long context becomes a second, separate bandwidth problem. Flag
  that, and say we return to it in section 3.

  Then the line that sets up the next slide: if the hardware wants hundreds of
  multiply-adds per byte and decode gives it one, most of the GPU is sitting
  idle, and we can say exactly how much.}
\end{frame}

\begin{frame}[shrink=6]{The roofline}
  \centering
  \begin{tikzpicture}
    \begin{loglogaxis}[
      width=8.2cm, height=5.0cm,
      xlabel={arithmetic intensity (FLOP per byte)},
      ylabel={achievable (TFLOP/s)},
      xmin=0.5, xmax=20000, ymin=1, ymax=3000,
      xtick={1,10,100,1000,10000},
      ytick={1,10,100,1000},
      grid=both, grid style={UHGrey!20},
      tick label style={font=\tiny}, label style={font=\scriptsize},
    ]
      \addplot[UHBlue, very thick] coordinates {(0.5,2) (247,990)};
      \addplot[UHGreen, very thick] coordinates {(247,990) (20000,990)};
      \addplot[only marks, mark=*, UHOrange, mark size=2.2pt] coordinates
        {(1,4) (32,128) (4096,990)};
      \draw[UHGrey, dashed] (axis cs:247,1) -- (axis cs:247,990);
      \node[lbl, anchor=west] at (axis cs:1.4,4.4) {decode, 1 user};
      \node[lbl, anchor=west] at (axis cs:45,115) {decode, 32 users};
      \node[lbl, anchor=east] at (axis cs:3200,640) {prefill};
      \node[lbl, anchor=west] at (axis cs:300,2.6) {corner at 247};
      \node[lbl, anchor=west, text=UHBlue] at (axis cs:0.7,40)
        {memory roof: 4.0\,TB/s};
      \node[lbl, anchor=east, text=UHGreen] at (axis cs:18000,1450)
        {compute roof: 990\,TFLOP/s bf16};
    \end{loglogaxis}
  \end{tikzpicture}

  \vspace{0.2em}
  \raggedright\small
  A \key{roofline} is a ceiling chart: whatever your code does, it cannot go
  above either line. The sloped part is the limit set by how fast memory feeds
  the chip. The flat part is where the multipliers themselves become the
  limit. This card can do about 990 trillion floating-point operations a second but fetch only
  4 terabytes a second, so it is not even nominally busy until you hand it
  about \textbf{250 FLOPs per byte}, which works out at \key{roughly
  250 tokens in every decode step.}
  \source{Roofline model: Williams, Waterman \& Patterson, CACM 2009.
    Hardware figures: NVIDIA GH200 (H100-class, 96\,GiB HBM3) product
    specifications.}
  \note{This is the central picture of the lecture. Spend time on it.

  Read it as a budget. The diagonal says ``this is all you can do if you are
  memory-limited'', the flat line says ``this is all you can do, full stop'',
  and any real kernel sits under both. The corner (the ``ridge point'') is
  $990 / 4.0 \approx 247$ FLOP per byte. Below that you are limited by
  bandwidth and above it by arithmetic.

  Then the sentence that should make them uncomfortable: at batch size 1,
  decode runs at about $4/990$, i.e. under half a percent of the GPU's
  arithmetic capability. We bought a 990 TFLOP/s machine and are using four of
  them. No software bug is involved. That is simply what generating one token
  at a time costs.

  Honest caveats: 990 TFLOP/s is dense bf16 peak, which no kernel achieves, and
  the ridge point shifts with precision. In FP8 the compute roof roughly
  doubles, so the ridge moves right and batching matters even more.}
\end{frame}

\begin{frame}{A bandwidth estimate for decode speed}
  \[
    \text{tokens per second} \;\lesssim\;
    \frac{\text{memory bandwidth}}{\text{bytes of weights}}
  \]
  \emph{In words:} how many times per second can the card read the whole model?
  \vspace{0.2em}
  \begin{columns}[T]
    \begin{column}{0.62\textwidth}
      \textbf{Your project model, 4B in bf16, on one GH200:}
      \begin{itemize}\small
        \item 4 billion weights, 2 bytes each: \textbf{8\,GB} to read for
          every single token
        \item the card moves \textbf{4.0\,TB} of memory per second
        \item 8\,GB goes into 4.0\,TB about \key{500 times}, so 500 tokens a
          second, i.e.\ 2\,ms each
      \end{itemize}
      \vspace{0.2em}
      That is the \emph{ceiling} for a single stream, however good the kernels
      are.
    \end{column}
    \begin{column}{0.38\textwidth}
      \textbf{Include the KV cache}
      \vspace{0.2em}
      \small
      At 8k context, this example also reads about 1.2\,GB of cached keys and
      values each step. The ceiling becomes
      \[
        4.0\,\mathrm{TB/s} \div 9.2\,\mathrm{GB}
        \approx \textbf{430 tokens/s}.
      \]
    \end{column}
  \end{columns}
  \source{NVIDIA GH200 Grace Hopper Superchip datasheet (HBM3 and HBM3e
    variants); CSC, \texttt{docs.csc.fi/computing/systems-roihu}.}
  \note{Make them do this one out loud with you. It is two numbers and one
  division, and they will use it more often than anything else today.

  In symbols, if they want it written down: bytes per token
  $\approx 4\times10^{9}\times 2 = 8\times10^{9}$, bandwidth
  $\approx 4.0\times10^{12}$ B/s, so tokens/s
  $\lesssim 4.0\times10^{12}/8\times10^{9} = 500$; with the cache,
  $4.0\times10^{12}/9.2\times10^{9} \approx 430$.

  Insist on the sourcing habit: the bandwidth figure is a spec-sheet number for
  a specific SKU, and the 96\,GiB and 144\,GB GH200 parts differ by about a
  fifth. Someone quoting ``GH200 bandwidth'' without saying which has already
  made a larger error than most of the optimisations we discuss today will fix.

  The KV-cache correction in the right column is where the cache first appears
  as a \emph{bandwidth} problem as well as a memory one: you re-read the whole
  cache every single step.}
\end{frame}

\begin{frame}{From the ceiling to a measured speed}
  \centering
  \begin{tabular}{@{}L{5.0cm} R{2.2cm} L{4.0cm}@{}}
    \toprule
    & \textbf{tokens/s} & \textbf{why lower}\\
    \midrule
    Paper ceiling, weights only & 500 & ---\\
    Ceiling including an 8k KV cache & $\approx$ 430 & you re-read the cache every step\\
    Realistic achieved bandwidth ($\approx$ 70\%) & $\approx$ 300 & imperfect access patterns\\
    What a server actually reports, $B=1$ & \textbf{150--250} & scheduler, Python, HTTP, sampling\\
    \bottomrule
  \end{tabular}

  \vspace{0.3em}
  \raggedright
  Measure single-stream TPOT, invert it, and compare the result with the
  estimate. A result within a factor of two is useful agreement for this
  simplified model. Then repeat at concurrency 8 and 32: aggregate throughput
  should rise while each request receives a smaller share of the server.
  \note{The honest framing: a back-of-envelope estimate that lands within
  $2\times$ of reality is a success. Engineers use these to decide whether an
  idea is worth trying at all.

  Tell them what ``something is wrong'' looks like in practice: measuring with
  the server still warming up, counting the prefill time inside TPOT, running
  on CPU by accident, or an output length so short that start-up dominates.
  The lab briefing on Thursday has the protocol.

  This slide is the course's recurring promise in miniature: every lecture
  contains at least one number they reproduce themselves.}
\end{frame}

\begin{frame}[shrink=4]{TTFT, TPOT, and who feels which}
  \[
    \text{total wait} \;=\;
    \underbrace{\mathrm{TTFT}}_{\text{reading the prompt}} \;+\;
    \underbrace{\mathrm{TPOT} \times \text{(tokens after the first)}}
               _{\text{generating the answer}}
  \]
  \emph{In words:} you wait once for the prompt, then once per token produced.
  \begin{columns}[T]
    \begin{column}{0.52\textwidth}
      \textbf{Worked, 4B model, 2{,}000-token prompt, 500 output tokens:}
      \small
      \begin{itemize}
        \item prefill: pushing 2{,}000 tokens through a 4B model is about
          16 trillion floating-point operations
        \item the card really achieves about 400 trillion a second, so:
          \key{40\,ms}
        \item decode: 500 tokens at 5\,ms each is \key{2.5\,s}
      \end{itemize}
      \vspace{0.4em}
      Prefill is \textbf{1.6\%} of the wall clock. Two thousand prompt tokens
      cost about as much time as \textbf{eight} generated ones.
    \end{column}
    \begin{column}{0.48\textwidth}
      \begin{block}{Who feels what}
        \small
        \textbf{TTFT} --- chat users, agents doing many short calls,
        anything with a spinner.\\[0.3em]
        \textbf{TPOT} --- anyone reading a streamed response. Anything under
        about 30\,ms per token feels instant, and by 100\,ms it starts to feel
        broken.\\[0.3em]
        \textbf{Neither} --- batch jobs. For offline scoring only
        \emph{throughput} matters, and you should batch as hard as the memory
        allows.
      \end{block}
    \end{column}
  \end{columns}
  \note{The prefill arithmetic in symbols, for the board: FLOPs
  $\approx 2Nn = 2\times4\times10^{9}\times 2000 = 1.6\times10^{13}$, divided
  by an achieved $4\times10^{14}$ FLOP/s gives 40\,ms. The full latency
  expression is $\mathrm{TTFT} + \mathrm{TPOT}\times(n_{\text{out}}-1)$.

  The 1.6\% figure usually surprises people, and it is the practical
  justification for prefix caching and for chunked prefill later on. Prefill
  takes very little wall-clock time, but it arrives in one lump, which is why
  it can stall everyone else.

  Reading speed as the TPOT yardstick is worth giving: a fast reader does
  roughly 5 words per second, which is maybe 7 tokens per second, i.e. 140\,ms
  per token. Any TPOT below about 30\,ms is already faster than anyone reads,
  so further latency work on a chat product buys nothing a user can perceive.
  Spend the effort on throughput.

  Third bullet is the one students forget: for their own evaluation runs, they
  are a batch job, and they should stop optimising latency entirely.}
\end{frame}

\begin{frame}{Which technique helps which phase}
  \centering
  \begin{tabular}{@{}L{4.4cm} C{2.4cm} C{2.4cm}@{}}
    \toprule
    \textbf{Technique} & \textbf{helps prefill} & \textbf{helps decode}\\
    \midrule
    Bigger batches & no (already busy) & \textbf{yes, hugely}\\
    INT4 weight-only quantisation & \warn{can hurt} & \textbf{yes}\\
    FP8 weights + activations & \textbf{yes} & \textbf{yes}\\
    KV-cache quantisation & slightly & \textbf{yes}\\
    Prefix caching & \textbf{yes, enormously} & no\\
    Speculative decoding & no & \textbf{yes}\\
    \bottomrule
  \end{tabular}
  \vfill
  \begin{takeaway}
    Work out which half of the request you are paying for before you optimise.
  \end{takeaway}
  \note{Use this table as the map for the rest of the lecture, since every row
  is a section heading later on. Promise them that and move on, and do not
  explain the rows now. Add out loud that decode is usually the half you are
  paying for, and that memory bandwidth is what limits it.

  Three more rows that belong on this map and that we will reach later:
  FlashAttention helps prefill a lot and decode somewhat; chunked prefill does
  not speed prefill up but stops it from blocking decode; and more GPUs in
  tensor parallel help prefill, and help decode too if the interconnect is
  fast enough.

  The one row to point at is INT4. It is a genuine trade rather than a free
  win, and it is one of today's three gotchas. Tell them it is coming in
  section 5 so they are listening for it.

  Section hand-off: we have established that decode moves eight gigabytes of
  weights per token. The next section is about the \emph{other} thing you have
  to move and store, which is bigger and which nobody budgets for.}
\end{frame}

% =====================================================================
\section{The KV cache}
% =====================================================================

\begin{frame}{Week 1's formula, read as a budget}
  \small
  \[
    \text{KV bytes} = \text{requests} \times \text{tokens} \times
    \text{layers} \times \text{KV heads} \times
    \underbrace{\text{numbers per head}}_{128} \times
    \underbrace{2}_{\text{key and value}} \times
    \underbrace{2}_{\text{bytes each, bf16}}
  \]
  \normalsize
  \emph{In words:} every layer stores two small vectors for every token of
  every request, and you pay for all of them at once.
  Your project model has 36 layers and 8 KV heads, so
  \[
    36 \times 8 \times 128 \times 2 \times 2 \;\approx\;
    147\unit{kB}\ \text{per token}
  \]
  \vspace{-0.3em}
  \begin{columns}[T]
    \begin{column}{0.5\textwidth}
      \begin{block}{In week 1 this was an architecture slide}
        \small
        It explained why GQA, MLA, SWA and cross-layer sharing exist.
      \end{block}
    \end{column}
    \begin{column}{0.5\textwidth}
      \begin{block}{Today it is a capacity planning slide}
        \small
        It tells you how many people you can serve, and it is the number your
        server runs out of first.
      \end{block}
    \end{column}
  \end{columns}
  \vspace{0.4em}
  \centering
  \small
  And this is not only a storage cost. Every decode step \emph{reads the whole
  cache back}, so it eats bandwidth too.
  \note{Clock check: should be around 16:38.

  The symbolic version, as it appeared in week 1:
  $B \times n \times L \times h_{kv} \times d_{\text{head}} \times 2 \times b$,
  with $B$ requests, $n$ tokens, $L=36$ layers, $h_{kv}=8$ KV heads,
  $d_{\text{head}}=128$ and $b=2$ bytes for bf16.

  Put the week 1 number on screen unchanged --- 147\,kB per token --- and say
  explicitly that it is the same slide with a different question attached.
  Repetition across lectures is deliberate; this is the third time they have
  seen this formula and it should now feel like arithmetic rather than theory.

  The closing line is the bit that is new today. In week 1 the cache was a
  memory cost. In serving it is a memory cost that caps concurrency and a
  bandwidth cost that slows every step, and you pay both.}
\end{frame}

\begin{frame}{The same model at 8k and at 128k}
  \centering
  \begin{tabular}{@{}L{3.2cm} R{2.6cm} R{2.6cm} R{2.4cm}@{}}
    \toprule
    \textbf{Context} & \textbf{KV / request} & \textbf{$\times$ weights (8\,GB)} & \textbf{of 96\,GiB}\\
    \midrule
    1 request @ 8k    & 1.2\,GB  & $0.15\times$ & 1\%\\
    1 request @ 32k   & 4.8\,GB  & $0.6\times$  & 5\%\\
    1 request @ 128k  & 19.3\,GB & $2.4\times$  & 19\%\\
    \bottomrule
  \end{tabular}

  \vspace{1.0em}
  \raggedright
  \begin{columns}[T]
    \begin{column}{0.52\textwidth}
      \textbf{One} user at 128k context occupies \key{more than twice the
      memory of the entire model}.
      \vspace{0.5em}

      The relationship is linear: double the context and you double the cache,
      which halves the number of users you can hold. There is no economy of
      scale in that formula.
    \end{column}
    \begin{column}{0.48\textwidth}
      \begin{block}{Units, stated once}
        \small
        96\,GiB $= 103$\,GB.\\
        A server typically reserves $\approx 90\%$ of the card, and the weights
        take 8\,GB, so
        \[
          0.9 \times 103 - 8 \;\approx\; \mathbf{85\unit{GB}}
        \]
        is what is actually available for KV. We use 85\,GB from here on.
      \end{block}
    \end{column}
  \end{columns}
  \note{The $2.4\times$ is the number that should stick. People budget GPU
  memory by looking at the model size, which is exactly the wrong thing to
  look at once contexts get long.

  Do the GiB-versus-GB correction openly. It is a 7\% difference, and it makes
  a capacity estimate wrong in the direction that hurts.

  The 85\,GB working figure carries through the next two slides, so write it
  on the board.}
\end{frame}

\begin{frame}[shrink=7]{Multiply by concurrency until it breaks}
  \centering
  \begin{tabular}{@{}R{2.0cm} R{2.4cm} R{2.6cm} L{3.6cm}@{}}
    \toprule
    \textbf{Concurrent} & \textbf{KV @ 8k} & \textbf{+ 8\,GB weights} & \\
    \midrule
    1   & 1.2\,GB  & 9.2\,GB  & trivially fine\\
    8   & 9.6\,GB  & 17.6\,GB & fine\\
    32  & 38.6\,GB & 46.6\,GB & fine --- \cit{week 1's number}\\
    64  & 77\,GB   & 85\,GB   & \warn{at the limit}\\
    70  & 84\,GB   & 92\,GB   & \bad{over budget}\\
    \bottomrule
  \end{tabular}

  \vspace{0.8em}
  \raggedright
  At \textbf{128k} the same arithmetic gives \textbf{4} concurrent requests:
  sixteen times the context buys you sixteen times fewer users. The card still
  has arithmetic to spare and nowhere left to put anyone.
  \vfill
  \begin{block}{The conclusion of this section}
    \key{In this example, the KV cache sets the concurrency limit.}
    The 4B model's weights take 8\,GB, while its cache can grow to many times
    that at long contexts and high concurrency.
  \end{block}
  \note{Walk down the table a row at a time and let them see the collision
  coming. The 32-request row is the figure from week 1, so they should
  recognise it; the point of the table is that week 1 stopped one row too
  early.

  The sentence in the block is takeaway number two of the lecture. Say it
  twice.

  Anticipate the obvious question --- ``so why not just buy more memory?'' ---
  and answer it now: memory is the scarce and expensive part of a GPU, and the
  cache grows linearly with both users and context, so you lose that race. The
  rest of this section is the alternatives.}
\end{frame}

\begin{frame}[shrink=6]{Where the memory actually goes under naive allocation}
  \centering
  \begin{tikzpicture}[x=1mm, y=1mm]
    \node[lbl, anchor=west] at (-2,26) {reserve \texttt{max\_model\_len} per request:};
    % request 1
    \fill[UHBlue!40, draw=UHBlue] (0,16) rectangle (18,22);
    \fill[UHGrey!15, draw=UHGrey!50] (18,16) rectangle (95,22);
    \node[lbl, anchor=east] at (-2,19) {req 1};
    \node[font=\tiny] at (9,19) {used};
    \node[font=\tiny, text=UHGrey] at (56,19) {reserved, never used};
    % request 2
    \fill[UHBlue!40, draw=UHBlue] (0,8) rectangle (40,14);
    \fill[UHGrey!15, draw=UHGrey!50] (40,8) rectangle (95,14);
    \node[lbl, anchor=east] at (-2,11) {req 2};
    \node[font=\tiny] at (20,11) {used};
    % request 3
    \fill[UHBlue!40, draw=UHBlue] (0,0) rectangle (9,6);
    \fill[UHGrey!15, draw=UHGrey!50] (9,0) rectangle (95,6);
    \node[lbl, anchor=east] at (-2,3) {req 3};
    \draw[decorate, decoration={brace, amplitude=4pt}, UHOrange, thick]
      (95,-2) -- (0,-2);
    \node[lbl, anchor=north, text=UHOrange] at (47,-7)
      {one slot sized for the worst case};
  \end{tikzpicture}

  \vspace{1.2em}
  \raggedright\small
  Three kinds of waste, all of them avoidable:
  \begin{itemize}
    \item \key{internal} --- you reserved 8k and the answer was 200 tokens
    \item \key{reservation} --- memory held for tokens that may never be generated
    \item \key{external fragmentation} --- free memory exists but not in one
      contiguous run of the right size
  \end{itemize}
  Measured waste in pre-paged serving systems: \textbf{60--80\%} of the KV
  region.
  \source{Kwon et al., \textit{Efficient Memory Management for LLM Serving with
    PagedAttention}, SOSP 2023.}
  \note{The analogy to use is memory allocation before virtual memory: if every
  process had to be given one contiguous physical region sized for its maximum
  possible footprint, you would waste most of your RAM and still fail to start
  the fourth process. That is exactly the state LLM serving was in until 2023.

  The 60--80\% figure comes from the vLLM paper and is worth quoting
  precisely. It means that one of the largest serving wins of recent years came
  from memory management rather than from a better model or a faster kernel.

  That sets up the next slide.}
\end{frame}

\begin{frame}{PagedAttention}
  \centering
  \begin{tikzpicture}[x=1mm, y=1mm, font=\scriptsize]
    % logical blocks A
    \node[lbl, anchor=west] at (0,30) {Sequence A, logical blocks};
    \foreach \i/\lab in {0/0,1/1,2/2,3/3}
      \node[tokcell, minimum width=8mm, fill=UHBlue!25] (A\i) at ({4+\i*10},24) {\lab};
    % logical blocks B
    \node[lbl, anchor=west] at (0,14) {Sequence B, logical blocks};
    \foreach \i/\lab in {0/0,1/1}
      \node[tokcell, minimum width=8mm, fill=UHOrange!25] (B\i) at ({4+\i*10},8) {\lab};
    % physical pool
    \node[lbl, anchor=west] at (62,30) {Physical KV blocks (one pool)};
    \foreach \i in {0,...,3}
      \node[tokcell, minimum width=8mm, fill=UHLight] (P\i) at ({66+\i*10},24) {P\i};
    \foreach \i/\n in {0/4,1/5,2/6,3/7}
      \node[tokcell, minimum width=8mm, fill=UHLight] (P\n) at ({66+\i*10},14) {P\n};
    % arrows
    \draw[flow] (A0) -- (P2);
    \draw[flow] (A1) -- (P5);
    \draw[flow] (A2) -- (P0);
    \draw[flow] (A3) -- (P7);
    \draw[flowhi] (B0) -- (P2);
    \draw[flowhi] (B1) -- (P4);
    \node[lbl, anchor=north, text=UHOrange] at (86,1)
      {P2 shared, refcount 2 --- copied only on write};
  \end{tikzpicture}

  \vspace{0.6em}
  \raggedright\small
  \begin{itemize}
    \item KV is stored in fixed-size \key{blocks} (16 tokens is typical),
      allocated on demand
    \item a per-sequence \key{block table} maps logical to physical blocks ---
      they need not be contiguous
    \item waste is now at most one partly-filled block per sequence:
      \key{under 4\%}
    \item identical prefixes can \key{share} a physical block; a block is
      copied only when one sequence writes to it
  \end{itemize}
  \source{Kwon et al., SOSP 2023 --- the paper vLLM is built on.}
  \note{Say the operating-systems word out loud: this is paging, applied to the
  KV cache, and the block table is a page table. Several of them have taken an
  OS course and the analogy does all the work.

  Copy-on-write matters more than it looks. Sampling $n=8$ completions from one
  prompt means eight sequences sharing every prompt block until they diverge,
  so best-of-$n$ costs far less memory than eight independent requests. That is
  directly relevant to anyone doing self-consistency in their project.

  The 4\% versus 60--80\% comparison is the headline. One idea, borrowed from
  1960s systems work, bought roughly a $2$--$4\times$ increase in serving
  throughput.}
\end{frame}

\begin{frame}{Prefix caching: paying for the system prompt once}
  \begin{columns}[T]
    \begin{column}{0.54\textwidth}
      If a prompt \key{begins with} tokens whose KV blocks are still cached,
      the server reuses them and skips that part of prefill.
      \vspace{0.5em}

      Systems such as SGLang and vLLM can find and reuse a shared prefix across
      many requests.
      \vspace{0.5em}

      This does not intentionally approximate the model, but it consumes cache
      memory. The gain depends on whether the prefix remains resident.
    \end{column}
    \begin{column}{0.46\textwidth}
      \begin{block}{Worked: a 5-shot evaluation run}
        \small
        Shared prefix: system prompt + 5 examples $\approx$ 1{,}200 tokens.\\
        Per-item suffix: $\approx$ 80 tokens.\\[0.4em]
        Without prefix caching you prefill 1{,}280 tokens per item.
        With it, \textbf{80}.\\[0.4em]
        Over a 1{,}000-item test set that is
        up to \key{$16\times$ less prefill work} while that prefix remains
        cached.
      \end{block}
      \vspace{0.2em}
      \small
      The same idea helps multi-turn chat, agent loops and fixed RAG
      instructions.
    \end{column}
  \end{columns}
  \source{Zheng et al., \textit{SGLang: Efficient Execution of Structured
    Language Model Programs}, 2024 (RadixAttention); vLLM automatic prefix
    caching.}
  \note{This is often a cheap win for evaluation workloads with a long shared
  prefix. The gain depends on cache residency and available memory.

  The ordering consequence is the practical tip. Put the stable part of your
  prompt \emph{first} and the variable part last. A template that interpolates
  the question into the middle of the system prompt destroys the shared prefix,
  and nobody notices, because nothing breaks; the server just gets slower.

  Forward pointer: in Dmitry's RAG lab the retrieved passages change every
  query, so the cacheable prefix is only the instruction block. Worth saying
  now so they design the prompt correctly in week 5.}
\end{frame}

% Architecture-specific serving examples were introduced in Lecture 1.
\iffalse
\begin{frame}{The week 1 attention variants, priced in users}
  \centering
  \small
  \begin{tabular}{@{}L{2.4cm} L{2.8cm} L{6.0cm}@{}}
    \toprule
    \textbf{Move} & \textbf{Mechanism} & \textbf{What it buys at serving time}\\
    \midrule
    \key{Share heads} & MQA, GQA &
      query heads share one key/value head; the cache shrinks by that factor\\
    \key{Compress} & MLA, CCA &
      store a squeezed-down version; the extra matrix multiply lands in a
      phase with compute to spare\\
    \key{Shorten} & SWA, DSA &
      most layers keep only a recent window, so the cache grows with the
      window rather than the context\\
    \key{Share layers} & cross-layer KV &
      later layers reuse earlier layers' keys and values\\
    \bottomrule
  \end{tabular}

  \vspace{0.5em}
  \raggedright
  In week 1 these were four ways to make a config file look better. Each one
  is really a line in the capacity table three slides ago: \key{more users on
  the same card}.
  \source{Raschka, \textit{A Visual Guide to Attention Variants}, Mar 2026;
    \textit{Recent Developments in LLM Architectures}, May 2026.}
  \note{Deliberate repetition of the week 1 taxonomy slide, with a different
  final column. Do not re-teach the mechanisms. Ask the room what each one does
  and let them answer.

  Two details trimmed off the slide and worth saying: a sliding window in DSA
  is 2{,}048 tokens, and cross-layer KV keeps each layer's own queries while
  reusing the earlier layers' keys and values.

  The precise statements, if pressed: GQA divides the cache by $h_q/h_{kv}$;
  MLA caches a low-rank latent and projects it back up; sliding-window and
  sparse attention take the cache from $O(L^2)$ to $O(L\cdot k)$ with
  $k=2048$ in DeepSeek's DSA; cross-layer KV reuses earlier layers' $K$ and
  $V$ tensors.

  The sentence at the bottom is the point of the slide and the through-line of
  the course: architecture decisions are cost decisions, and today is where
  the cost turns into a number of users.

  Keep this quick, ninety seconds. The next slide has the numbers.}
\end{frame}

\begin{frame}{What those moves are worth in gigabytes}
  \begin{columns}[T]
    \begin{column}{0.5\textwidth}
      \textbf{Cross-layer KV --- Gemma 4 E2B}\\[0.3em]
      \small
      35 layers; the first 15 compute their own $K,V$, the final 20 reuse
      them. Roughly half the cache disappears:
      \key{$\approx$ 2.7\,GB saved at 128k context in bf16} (E4B: $\approx$
      6\,GB).
      \vspace{0.8em}

      \textbf{Linear-attention hybrids --- Kimi Linear}\\[0.3em]
      \small
      Kimi Delta Attention at 3 linear : 1 full. Reported against full
      attention: \key{75\% KV cache reduction} and \key{up to $6\times$ decode
      throughput}.
      \vspace{0.4em}
      The mechanism is the one from week 1: a linear layer's state is a fixed
      size, \emph{independent of context length}.
    \end{column}
    \begin{column}{0.5\textwidth}
      \begin{block}{Read these claims carefully}
        \small
        ``Up to $6\times$'' is a vendor-reported figure at a favourable context
        length and batch size, against the vendor's own full-attention
        baseline. The measurement is probably sound. You should not expect it
        on your own workload.
        \vspace{0.4em}
        The 75\% cache figure is the more robust of the two, because it follows
        from the architecture rather than from a kernel.
      \end{block}
      \vspace{0.3em}
      \small
      Both of these are decode and concurrency wins. Neither changes prefill
      much.
    \end{column}
  \end{columns}
  \sources{Gemma 4 model card, as summarised in Raschka,
    \textit{Recent Developments in LLM Architectures}, May 2026;
    Kimi Linear technical report (48B), Oct 2025.}
  \note{Two concrete numbers, both dated, both attributed, one of them
  deliberately questioned. That pairing is the research-literacy lesson: quote
  the number, then say what kind of number it is.

  The distinction worth teaching: an architectural claim (``we store half as
  many KV tensors'') is checkable from the config file and cannot really be
  wrong. A throughput claim (``$6\times$ faster'') is a measurement on someone
  else's hardware with someone else's kernels and someone else's batch size.
  Trust them differently.

  That scepticism leads straight into the caveat on the next slide.}
\end{frame}
\fi

\begin{frame}{Capacity, summarised, with one caveat}
  \centering
  \begin{tabular}{@{}L{2.4cm} R{2.6cm} R{3.0cm} R{3.0cm}@{}}
    \toprule
    \textbf{Context} & \textbf{KV / request} & \textbf{max concurrent, bf16 KV} & \textbf{with FP8 KV}\\
    \midrule
    8k   & 1.2\,GB  & 70 & 140\\
    32k  & 4.8\,GB  & 17 & 35\\
    128k & 19.3\,GB & 4  & 8\\
    \bottomrule
  \end{tabular}

  \vspace{0.4em}
  \raggedright\small
  One GH200, 4B model in bf16, 85\,GB available for KV.
  \vspace{0.7em}
  \begin{block}{A practitioner's caveat, worth repeating all term}
    \small
    A model with a theoretically cheaper attention mechanism is not
    automatically cheaper to \emph{serve}. Inference stacks have years of
    optimisation behind plain GQA and far less behind hybrid and linear
    attention. In late 2026 a ``slower'' classic GQA model can still serve faster
    in practice than a hybrid one.
    \vspace{0.3em}
    Until the kernels catch up, \key{the better-supported method usually wins}.
  \end{block}
  \note{Clock check: should be around 16:52 as this comes down.

  The table is the deliverable of the section: three context lengths, three
  concurrency numbers, one quantisation column that previews section 5. Tell
  them to copy it, because the structure (not the values) is what they will
  reproduce for their own model.

  The caveat is a genuine field observation and it generalises beyond
  attention: DSA, MLA and linear attention all took time to get fast, and a
  published FLOP count and the wall-clock time on your own box are rarely the
  same thing. It is also the honest reason I would still default to a GQA model
  for a production deployment today.}
\end{frame}

% =====================================================================
\section{Continuous batching}
% =====================================================================

\begin{frame}[shrink=9]{Static batching wastes most of the GPU}
  \centering
  \begin{tikzpicture}[x=1mm, y=1mm, font=\tiny]
    \draw[UHGrey, ->] (0,-4) -- (104,-4) node[lbl, anchor=west] {time};
    \foreach \i/\len/\y in {1/30/24, 2/52/16, 3/70/8, 4/88/0} {
      \fill[UHBlue!35, draw=UHBlue] (0,\y) rectangle (\len,{\y+6});
      \fill[UHGrey!15, draw=UHGrey!40] (\len,\y) rectangle (88,{\y+6});
      \node[lbl, anchor=east] at (-2,{\y+3}) {req \i};
    }
    \node at (15,27) {};
    \node[font=\tiny, text=UHGrey] at (60,27) {idle slot --- GPU does the work anyway};
    \draw[UHOrange, thick, dashed] (88,-3) -- (88,32);
    \node[lbl, anchor=south, text=UHOrange] at (88,32) {next batch can start};
    % second batch
    \foreach \y in {24,16,8,0} {
      \fill[UHBlue!35, draw=UHBlue] (90,\y) rectangle (104,{\y+6});
    }
  \end{tikzpicture}

  \vspace{1.0em}
  \raggedright
  \begin{itemize}
    \item The batch is fixed when it starts and released when the
      \emph{slowest} member finishes.
    \item A finished sequence still occupies a slot; the kernel computes
      padding for it.
    \item Request 5 waits behind request 4's last token even though three
      slots are free.
  \end{itemize}
  \key{Output lengths vary by an order of magnitude} in real traffic. A
  one-line answer and a 4{,}000-token reasoning trace end up in the same
  batch, so this waste is the normal case.
  \source{Yu et al., \textit{Orca: A Distributed Serving System for
    Transformer-Based Generative Models}, OSDI 2022.}
  \note{Point at the grey. That is a GPU that costs two or three euros an hour,
  running at full power, computing nothing anyone asked for.

  Tie it back to section 1: reasoning effort is a per-request parameter, so the
  spread in output lengths is something your users control rather than an
  accident of the workload. That makes static batching worse over time.

  Historical note worth thirty seconds: this is how everyone served models
  until Orca in 2022. It is a very recent fix for a very obvious problem, which
  should encourage anyone in the room who thinks all the good systems ideas are
  taken.}
\end{frame}

\begin{frame}[shrink=9]{Continuous (iteration-level) batching}
  \centering
  \begin{tikzpicture}[x=1mm, y=1mm, font=\tiny]
    \draw[UHGrey, ->] (0,-4) -- (104,-4) node[lbl, anchor=west] {time};
    % req1 finishes at 30, req5 starts
    \fill[UHBlue!35, draw=UHBlue] (0,24) rectangle (30,30);
    \fill[UHGreen!30, draw=UHGreen] (31,24) rectangle (78,30);
    \node[lbl, anchor=east] at (-2,27) {slot 1};
    \node at (15,27) {req 1}; \node at (54,27) {req 5};

    \fill[UHBlue!35, draw=UHBlue] (0,16) rectangle (52,22);
    \fill[UHGreen!30, draw=UHGreen] (53,16) rectangle (90,22);
    \node[lbl, anchor=east] at (-2,19) {slot 2};
    \node at (26,19) {req 2}; \node at (71,19) {req 6};

    \fill[UHBlue!35, draw=UHBlue] (0,8) rectangle (70,14);
    \fill[UHGreen!30, draw=UHGreen] (71,8) rectangle (104,14);
    \node[lbl, anchor=east] at (-2,11) {slot 3};
    \node at (35,11) {req 3}; \node at (87,11) {req 7};

    \fill[UHBlue!35, draw=UHBlue] (0,0) rectangle (88,6);
    \fill[UHGreen!30, draw=UHGreen] (89,0) rectangle (104,6);
    \node[lbl, anchor=east] at (-2,3) {slot 4};
    \node at (44,3) {req 4};
  \end{tikzpicture}

  \vspace{1.0em}
  \raggedright
  The scheduler rebuilds the batch \key{every decoding iteration}, not every
  request:
  \begin{itemize}
    \item a sequence that emits its stop token leaves immediately and frees
      its KV blocks
    \item a waiting request joins at the very next step, prefill included
    \item padding, idle slots and head-of-line blocking all disappear
  \end{itemize}
  \vspace{0.3em}
  Reported throughput gains over static batching at the same latency:
  roughly \key{$2$--$4\times$}, before PagedAttention is even considered.
  \source{Yu et al., OSDI 2022; Kwon et al., SOSP 2023.}
  \note{The key word is \emph{iteration-level}. The unit of scheduling stops
  being ``a request'' and becomes ``one forward pass'', which is the right
  granularity because every decode step is independent of the last except
  through the KV cache.

  Note how this depends on the previous section: you can only admit a new
  sequence mid-flight if you can allocate its KV incrementally, which is
  exactly what paged blocks give you. Orca and PagedAttention are complements,
  and together they are most of why a modern server is an order of magnitude
  better than a naive loop.

  No green block starts before 31 on the diagram --- point that out, it is the
  ``join at the next step'' behaviour made visible.}
\end{frame}

\begin{frame}{Throughput and latency pull in opposite directions}
  \centering
  \begin{tikzpicture}
    \begin{axis}[
      width=7.8cm, height=4.8cm,
      xlabel={server throughput (output tokens/s)},
      ylabel={per-user TPOT (ms)},
      xmin=0, xmax=4000, ymin=0, ymax=140,
      xtick={0,1000,2000,3000,4000},
      ytick={0,25,50,75,100,125},
      tick label style={font=\tiny}, label style={font=\scriptsize},
      grid=both, grid style={UHGrey!20},
    ]
      \addplot[UHBlue, very thick, mark=*, mark size=1.3pt] coordinates
        {(200,9) (800,11) (1600,15) (2400,22) (3000,34) (3400,50)
         (3600,78) (3700,130)};
      \draw[UHRed, thick, dashed] (axis cs:0,50) -- (axis cs:4000,50);
      \node[lbl, anchor=west, text=UHRed] at (axis cs:60,57) {latency SLO: 50\,ms};
      \fill[UHGreen, opacity=0.10] (axis cs:0,0) rectangle (axis cs:3400,50);
      \node[lbl, anchor=south, text=UHGreen] at (axis cs:1700,2) {goodput region};
      \draw[UHGrey, dashed] (axis cs:3400,0) -- (axis cs:3400,50);
      \node[lbl, anchor=west] at (axis cs:3420,20) {max goodput};
    \end{axis}
  \end{tikzpicture}

  \vspace{0.2em}
  \raggedright\small
  Raise concurrency and aggregate throughput climbs, until the knee, where
  each added request slows everyone without adding tokens. \key{Goodput} is
  throughput that still meets the latency target, and past the knee extra load
  produces only violations.
  \source{Shape after published vLLM and DistServe throughput--latency sweeps;
    values illustrative.}
  \note{Make them read the axes carefully: \emph{more throughput is to the
  right, worse latency is up}. The curve going up and to the right means you
  buy throughput with latency, which is the trade the whole section is about.

  Two operating points worth naming. On the far left there is one user, lovely
  latency, and a GPU running at half a percent of peak. Near the knee you get a
  hundred times the throughput at a latency a human can still live with.

  The SLO line is where engineering meets product. There is no ``optimal''
  operating point until someone states a latency target, and if nobody does you
  end up tuned for a benchmark instead of for your users.

  Thursday they sweep concurrency at 1, 8 and 32 and draw this curve for their
  own model.}
\end{frame}

\begin{frame}{Chunked prefill}
  \begin{columns}[T]
    \begin{column}{0.52\textwidth}
      \textbf{The problem.} A 32k-token prompt arrives. Its prefill is one
      large, compute-bound forward pass. For its duration, every other user's
      decoding stops.
      \vspace{0.4em}

      One 32k prefill at $\approx$ 400\,TFLOP/s is roughly \key{650\,ms}, and
      that shows up as a 650\,ms stall in somebody else's token stream, for a
      request they have nothing to do with.
      \vspace{0.7em}

      \textbf{The fix.} Split the prefill into chunks (for example 512 tokens)
      and schedule each chunk \emph{alongside} the decode steps of the running
      batch.
    \end{column}
    \begin{column}{0.48\textwidth}
      \begin{block}{Why it also helps throughput}
        \small
        A decode-only step spends its time waiting on memory and leaves the
        arithmetic units idle. A prefill chunk is the opposite kind of work:
        plenty of arithmetic, hardly any new weights to fetch.
        \vspace{0.3em}
        \key{Mixing them in one batch uses both resources at once.}
      \end{block}
      \vspace{0.2em}
      \small
      The cost: TTFT for the long prompt gets slightly worse, because its
      prefill is now interleaved rather than monolithic. You are trading one
      user's TTFT for everyone else's TPOT, deliberately.
    \end{column}
  \end{columns}
  \source{Agrawal et al., \textit{Taming Throughput--Latency Tradeoff in LLM
    Inference with Sarathi-Serve}, OSDI 2024; implemented in vLLM and SGLang.}
  \note{The fairness framing is the memorable one: without chunked prefill, one
  user with a long prompt can freeze the service for everyone, and the person
  who suffers has no idea why.

  The complementarity argument in the right-hand block is the elegant part and
  worth drawing on the roofline from section 2: a mixed batch sits between the
  two extreme points, using some bandwidth and some arithmetic, rather than
  saturating one and wasting the other.

  Practical note for Thursday: chunked prefill is on by default in current
  vLLM. If their TTFT measurements look oddly high for long prompts under
  concurrency, this is why, and it is correct behaviour.}
\end{frame}

\begin{frame}[fragile]{The flags that matter}
\begin{lstlisting}[language=slurm]
# annotated; in the lab these are all on one line
vllm serve Qwen/Qwen3.5-4B-Instruct
  --max-model-len 8192           # caps KV per request: 1.2 GB here
  --max-num-seqs 64              # concurrency ceiling (slots)
  --max-num-batched-tokens 8192  # tokens per scheduler step = prefill chunk
  --gpu-memory-utilization 0.90  # fraction of the card vLLM may claim
  --enable-prefix-caching        # free, turn it on
  --enable-lora --lora-modules proj=./lora-week2
  --port 8000
\end{lstlisting}
  \vspace{0.3em}
  \small
  \begin{tabular}{@{}L{4.2cm} L{7.2cm}@{}}
    \texttt{max\_model\_len} & the biggest lever on concurrency; do not leave
      it at the model's 128k default if your task uses 2k\\
    \texttt{max\_num\_seqs} & raising it past what KV memory allows just moves
      the queue inside the server\\
    \texttt{max\_num\_batched\_tokens} & bigger $\Rightarrow$ better throughput,
      worse tail TPOT\\
    \texttt{gpu\_memory\_utilization} & on a shared GH200, lower it or your
      neighbour's job dies\\
  \end{tabular}
  \note{Walk the flags. The one people get wrong is
  \texttt{max\_model\_len}: it defaults to whatever the model config claims,
  which for a 2026 model is often 128k, and that single default can cut
  concurrency from seventy to four. It is the capacity table from section 3,
  as a command-line argument.

  \texttt{gpu\_memory\_utilization} has a social dimension on Roihu: vLLM
  pre-allocates its KV pool at start-up, so a server left running holds the
  memory whether or not anyone is using it. Repeat the compute-etiquette rule
  from week 1.

  Tell them this exact command, with their own model, is roughly what Thursday
  starts with.}
\end{frame}

\begin{frame}[shrink=3]{Goodput: throughput that met the deadline}
  \begin{columns}[T]
    \begin{column}{0.52\textwidth}
      \textbf{Throughput} --- tokens per second, summed over all users.
      Easy to maximise: raise concurrency until memory runs out.
      \vspace{0.5em}

      \textbf{Goodput} --- tokens per second \emph{from requests that met the
      latency target}. Everything else is counted as zero.
      \vspace{0.6em}

      Under a 50\,ms TPOT target, a server at 3{,}700 tokens/s and 130\,ms has
      \key{a throughput of 3{,}700 and a goodput of 0}.
    \end{column}
    \begin{column}{0.48\textwidth}
      \begin{block}{How to report honestly}
        \small
        \begin{itemize}
          \item warm up before timing
          \item fix the output length
          \item separate TTFT from TPOT
          \item sweep concurrency: 1 / 8 / 32
          \item report p50 \textbf{and} p95, not the mean
        \end{itemize}
      \end{block}
    \end{column}
  \end{columns}
  \vfill
  \begin{takeaway}
    Continuous batching fills GPU time that was being wasted. It gives you no
    extra memory, so the ceiling is still the KV cache, and someone still has
    to choose where on the curve you want to sit.
  \end{takeaway}
  \note{Clock check: should be around 17:02.

  The 3{,}700-and-zero example is deliberately absurd and it makes the
  definition stick. Every published throughput number should prompt the
  question ``at what latency?'', and most of them do not answer it.

  The reporting checklist is Thursday's protocol in miniature, so read it
  properly. It is what turns a lab report into evidence. The mean-versus-p95
  point especially: users experience the tail, and a mean hides it completely.}
\end{frame}

% =====================================================================
\section{Quantisation for inference}
% =====================================================================

\begin{frame}{Three independent places to use fewer bits}
  \centering
  \begin{tabular}{@{}L{2.4cm} L{3.6cm} L{5.2cm}@{}}
    \toprule
    \textbf{What} & \textbf{Size, 4B model} & \textbf{What lowering it buys}\\
    \midrule
    \key{Weights} & 8\,GB in bf16 &
      fewer bytes read per decode step $\Rightarrow$ directly faster decode\\[0.3em]
    \key{Activations} & small, transient &
      lets you use the FP8 tensor cores $\Rightarrow$ faster \emph{prefill}\\[0.3em]
    \key{KV cache} & 1.2\,GB per 8k request &
      more concurrent users, and less to re-read every step\\
    \bottomrule
  \end{tabular}

  \vspace{0.4em}
  \raggedright
  These are \key{orthogonal choices}. ``The model is quantised'' is as
  uninformative as ``the model uses sliding-window attention'' was in week 1.
  Always ask \emph{which of the three} uses fewer bits, and in what format.
  \note{Clock check: should be around 17:02--17:04 starting this section.

  The three-way split is the organising idea and it maps exactly onto the
  prefill/decode split from section 2: weights are a decode problem,
  activations are a prefill problem, KV is a concurrency problem. Say that
  mapping out loud; it makes the rest of the section predictable.

  The parallel with the SWA gotcha from week 1 is deliberate. The lesson
  transfers: a named technique without its parameters tells you very little.}
\end{frame}

\begin{frame}{FP8 on a GH200}
  \begin{columns}[T]
    \begin{column}{0.5\textwidth}
      \textbf{What it is}
      \begin{itemize}
        \item weights and activations use 8-bit floating-point numbers;
        \item each tensor or block also stores a scale, because FP8 represents
          a much narrower range than bf16;
        \item smaller blocks use more scales but preserve outliers better.
      \end{itemize}
      Hopper tensor cores multiply FP8 directly, so this is a native hardware
      format rather than compressed storage that must first be unpacked.
    \end{column}
    \begin{column}{0.5\textwidth}
      \textbf{Why it is used}
      \begin{itemize}
        \item weights halve: 8\,GB $\to$ 4\,GB;
        \item less weight traffic speeds up decode;
        \item faster FP8 tensor cores can also speed up prefill.
      \end{itemize}
      \textbf{How to evaluate it}
      Re-run the same held-out scorer and report quality beside latency and
      throughput. Published average losses do not guarantee that your
      language, task or long-context workload is unaffected.
    \end{column}
  \end{columns}
  \source{NVIDIA Hopper FP8 documentation; vLLM FP8 quantisation guide.}
  \note{Outliers are the whole story of activation quantisation, and worth
  thirty seconds: a handful of channels in a transformer carry values hundreds
  of times larger than the rest, and a single per-tensor scale must accommodate
  them, which crushes everything else to a few representable levels.
  Per-channel scaling is the standard answer.

  Emphasise ``native'': FP8 on this hardware is not emulated, so you get faster
  arithmetic as well as fewer bytes. That is why FP8 helps both phases while
  INT4 helps only one, which is the next slide.}
\end{frame}

\begin{frame}[shrink=3]{INT4 weight-only: AWQ and GPTQ}
  \begin{columns}[T]
    \begin{column}{0.5\textwidth}
      Store each weight in 4 bits, \key{unpack it back to 16-bit on the GPU
      just before use}, and do the arithmetic in bf16. Only the storage
      shrinks; the activations stay full precision.
      \vspace{0.6em}

      \textbf{GPTQ} --- goes layer by layer on the finished model, choosing how
      to round each weight so the layer's \emph{output} changes as little as
      possible on a sample of real data.
      \vspace{0.5em}

      \textbf{AWQ} --- observes that a small fraction of weight channels matter
      far more, identifies them from \emph{activation} magnitudes, and rescales
      so those channels survive quantisation.
    \end{column}
    \begin{column}{0.5\textwidth}
      \begin{block}{Why weight-only works at all}
        \small
        It helps for exactly the reason section 2 gave: decode is waiting on
        memory. Shrinking the weights from 8\,GB to 2\,GB means a quarter of
        the bytes per token, so the head-estimate ceiling goes from 500 to
        \key{2{,}000 tokens/s}.
        \vspace{0.3em}
        When decode is strongly bandwidth-bound, much of the unpacking cost
        can be hidden by arithmetic units that would otherwise be idle.
      \end{block}
      \vspace{0.2em}
      \small
      Both need \key{calibration data} --- a few hundred sequences. Use data
      that looks like your task. Calibrating a Finnish NER model on English
      web text is a real and common own-goal.
    \end{column}
  \end{columns}
  \sources{Frantar et al., \textit{GPTQ}, ICLR 2023;
    Lin et al., \textit{AWQ}, MLSys 2024.}
  \note{The ceiling arithmetic: $4.0\times10^{12} / 2\times10^{9} = 2000$
  tokens/s, i.e.\ the same division as the head-estimate slide with 2\,GB of
  weights instead of 8. GPTQ is one-shot post-training quantisation using
  second-order (Hessian) information about the layer's reconstruction error;
  the ``unpacking'' step is dequantisation inside the matmul kernel.

  The sentence to land: weight-only quantisation is a \emph{bandwidth}
  optimisation, and it pays off only because of where decode sits on the
  roofline. Students who understand that will predict the gotcha on the next
  slide before you show it.

  Calibration is the practical trap. The calibration set defines what ``matters''
  for AWQ and what error GPTQ minimises, so it should resemble the deployment
  distribution. For this course that means: if your project is Finnish, do not
  accept an off-the-shelf INT4 checkpoint calibrated on English without
  re-scoring it.}
\end{frame}

\begin{frame}{Why INT4 can slow down prefill}
  Weight-only INT4 can speed up decode and \textbf{slow down prefill}.
  Decode benefits from reading fewer weight bytes. Prefill is compute-bound,
  so unpacking 4-bit weights to bf16 adds work to the phase that is already
  short of arithmetic capacity. FP8 is multiplied directly by Hopper tensor
  cores and can therefore help both phases.
  \vspace{0.4em}
  \centering
  \begin{tabular}{@{}L{3.4cm} C{2.6cm} C{2.6cm} C{2.2cm}@{}}
    \toprule
    & \textbf{bytes/token} & \textbf{decode} & \textbf{prefill}\\
    \midrule
    bf16 (W16A16) & 8\,GB & baseline & baseline\\
    FP8 (W8A8)    & 4\,GB & $\approx 1.8\times$ & $\approx 1.5\times$\\
    INT4 (W4A16)  & 2\,GB & $\approx 3\times$ & \bad{$0.8$--$1.0\times$}\\
    \bottomrule
  \end{tabular}

  \vspace{0.3em}
  \raggedright\small
  Indicative only: the prefill number depends heavily on the kernel and the
  prompt length. Read the \emph{sign} rather than the magnitude.
  \source{Behaviour reported across vLLM, TensorRT-LLM and llama.cpp kernel
    benchmarks; see also Lin et al., AWQ, MLSys 2024.}
  \note{Gotcha one of three. Give it a full minute and make them explain it
  back to you, because it is a pure test of whether the roofline from section 2
  actually landed.

  The general principle underneath: an optimisation that removes a constraint
  you are not hitting, at the cost of a resource you are short of, is a
  pessimisation. That is a transferable idea well beyond LLMs.

  Practical consequence for a long-prompt workload such as RAG in week 5: if
  most of your tokens are prompt tokens, INT4 may be the wrong choice
  altogether, and FP8 is the safer default on this hardware.}
\end{frame}

\begin{frame}{Quantising the KV cache}
  \begin{columns}[T]
    \begin{column}{0.5\textwidth}
      The third, independent knob. Store the cached keys and values in 8-bit
      FP8 instead of 16-bit bf16, and every number in the cache halves:
      147\,kB per token becomes \key{73.5\,kB}.
      \small
      \begin{itemize}
        \item 8k context: 1.2\,GB $\to$ 0.6\,GB per request
        \item concurrent users at 8k: 70 $\to$ \key{140}
        \item concurrent users at 128k: 4 $\to$ \key{8}
        \item and half the bytes re-read every decode step
      \end{itemize}
      \vspace{0.4em}
      INT4 caches exist too, but need a finer-grained set of scale factors to
      hold up.
    \end{column}
    \begin{column}{0.5\textwidth}
      \begin{block}{Where the damage shows up}
        \small
        Long-context behaviour can be \textbf{more sensitive} to KV
        quantisation: early tokens' quantised keys may be attended to thousands
        of times.
        \vspace{0.3em}
        So the headline benchmark at 2k context looks untouched, and the
        128k-context retrieval task quietly degrades. \key{Short-context
        evaluation cannot detect this failure.}
      \end{block}
      \vspace{0.3em}
      \small
      Position encoding makes the keys harder to quantise than the values:
      rotating them spreads their values out.
    \end{column}
  \end{columns}
  \source{Hooper et al., \textit{KVQuant}, NeurIPS 2024; vLLM FP8 KV cache
    documentation.}
  \note{The detail behind ``finer-grained scale factors'': INT4 KV generally
  needs per-channel scaling for the keys $K$ and per-token scaling for the
  values $V$. And RoPE is what makes $K$ the harder of the two --- rotated keys
  have a wider, position-dependent distribution than values do.

  This knob buys the most memory for the least risk, and it is usually the last
  one people reach for, because it is not what ``quantisation'' means in most
  blog posts.

  The accumulation argument is the technical heart of the slide. Weight
  quantisation error is fixed and applied uniformly; KV error is written once
  per token and then read back for the rest of the generation, so its influence
  grows with context length.

  Which is a direct instruction for their evaluation: if they quantise the KV
  cache, they must evaluate at the context length they intend to deploy at, not
  at the length that is convenient.}
\end{frame}

\begin{frame}{The quality cost, and where it shows up}
  \begin{columns}[T]
    \begin{column}{0.52\textwidth}
      For FP8 and good INT4 the published quality loss is usually
      \textbf{under a point} on MMLU-style benchmarks. That average hides
      where the damage actually lands:
      \begin{itemize}\small
        \item \key{long-tail inputs}: rare formats, unusual prompts
        \item \key{non-English text}, especially morphologically rich languages
        \item \key{long reasoning chains}, where one bad token derails the rest
        \item \key{long contexts}, if you quantised the KV cache
      \end{itemize}
    \end{column}
    \begin{column}{0.48\textwidth}
      \begin{block}{Which is why Thursday makes you score it}
        \small
        Serve the model in FP8, run the same held-out scorer used in week 2,
        and put the score in the same table as the tokens per second.
        \vspace{0.4em}
        \key{Always report the quality number next to the speed number.}
      \end{block}
    \end{column}
  \end{columns}
  \vfill
  \begin{takeaway}
    FP8 is the safe default on a GH200 and INT4 is a bet on decode. Either
    way, the only quality number worth anything is the one you measured on
    your own task.
  \end{takeaway}
  \note{Clock check: should be around 17:11.

  Two things trimmed from the slide. First, weights, activations and the KV
  cache are three separate decisions, and you can take them independently.
  Second, the non-English bullet is really about tokenizer fertility: a
  morphologically rich language needs more tokens per word, so quantisation
  error has more steps to accumulate over.

  The multilingual point is not a throwaway. A benchmark average over mostly
  English, mostly short, mostly well-formed inputs is the measurement least
  able to see quantisation damage, and several of them have Finnish projects
  where this is a live problem rather than a hypothetical one.

  They already have a task scorer from week 2 and a held-out set from week 1,
  so the extra work here is small. The last line of the block is the course's
  recurring demand in one sentence. Say it, then move on: the next
  section is the only technique today that is genuinely free.}
\end{frame}

% =====================================================================
\section{Speculative and constrained decoding}
% =====================================================================

\begin{frame}{Speculative decoding}
  \centering
  \begin{tikzpicture}[x=1mm, y=1mm, font=\scriptsize]
    % drafter
    \node[lbl, anchor=east] at (-2,30) {small drafter};
    \foreach \i/\t in {0/{the}, 1/{cat}, 2/{sat}, 3/{under}} {
      \node[tokcell, minimum width=13mm, fill=UHLight] (d\i) at ({8+\i*16},30) {\t};
    }
    \node[lbl, anchor=west] at (72,30) {$k=4$ proposals, 4 cheap passes};
    % target
    \fill[UHBlue!20, draw=UHBlue] (0,16) rectangle (66,24);
    \node at (33,20) {target model: \textbf{one} forward pass over all 4+1 positions};
    \node[lbl, anchor=east] at (-2,20) {verify};
    % result
    \node[lbl, anchor=east] at (-2,6) {result};
    \node[tokon, minimum width=13mm] at (8,6) {the};
    \node[tokon, minimum width=13mm] at (24,6) {cat};
    \node[tokon, minimum width=13mm] at (40,6) {sat};
    \node[tokmask, minimum width=13mm] at (56,6) {under};
    \node[tokcell, minimum width=13mm, fill=UHOrange!30] at (72,6) {on};
    \node[lbl, anchor=west] at (80,6) {resampled by the target};
    \draw[flow] (d0) -- (8,24.5);
    \draw[flow] (d3) -- (56,24.5);
    \node[lbl, anchor=north, text=UHGreen] at (24,1) {3 accepted};
    \node[lbl, anchor=north, text=UHRed] at (58,1) {1 rejected};
  \end{tikzpicture}

  \vspace{1.2em}
  \raggedright
  Four target-model tokens' worth of progress from \key{one} target forward
  pass. The target pass costs almost the same as a single-token pass, because
  verifying 5 positions is still a tiny batch on a bandwidth-bound GPU, and
  the weights are read once either way.
  \sources{Leviathan et al., \textit{Fast Inference from Transformers via
    Speculative Decoding}, ICML 2023;
    Chen et al., \textit{Accelerating LLM Decoding with Speculative Sampling},
    2023.}
  \note{The intuition to give before the mechanism: most tokens in most text
  are easy. ``The cat sat on the ...'' does not need a 4B model. A 0.5B drafter
  gets the easy ones right, and the big model's job becomes checking rather
  than generating.

  The reason it is a win at all is the roofline again. Running the target over
  five positions instead of one moves you slightly right along the bandwidth
  roof at essentially no extra wall-clock cost, because you read the same eight
  gigabytes of weights. Speculative decoding is a way to buy arithmetic
  intensity you were not using.

  Next slide: why this is exact and not an approximation. That is the part
  people do not believe.}
\end{frame}

\begin{frame}[shrink=4]{Why the output distribution is unchanged}
  The drafter proposes a token. The big model already computed how likely
  \emph{it} thinks that token was. Compare the two numbers:
  \[
    \text{accept with probability }
    \min\!\left(1,\;
      \frac{\text{big model's probability of the token}}
           {\text{drafter's probability of the token}}\right)
  \]
  \emph{In words:} keep the guess outright if the big model likes it at least
  as much as the drafter did; otherwise keep it only sometimes. If it is
  rejected, draw the replacement from the \key{residual} --- the part of the
  big model's distribution the drafter had been under-weighting.
  \vspace{0.2em}
  \begin{columns}[T]
    \begin{column}{0.56\textwidth}
      \begin{block}{The guarantee}
        \small
        Under the algorithm's assumptions, the tokens are distributed
        \textbf{exactly} as if they had been sampled from the big model alone.
        \vspace{0.3em}
        Not approximately, and not ``close enough in practice''. The procedure
        is ordinary rejection sampling, so the result is identical by
        construction.
      \end{block}
    \end{column}
    \begin{column}{0.44\textwidth}
      \small
      Consequences:
      \begin{itemize}
        \item a poor drafter mainly reduces the potential speed-up
        \item you still measure latency and run output regression tests
        \item the drafter can be anything: a small model, an $n$-gram table,
          the target's own extra heads
      \end{itemize}
    \end{column}
  \end{columns}
  \note{Write the formulas out if the room wants them. With $q$ the drafter's
  distribution and $p$ the target's at the same position, a proposal
  $x \sim q$ is accepted with probability $\min(1, p(x)/q(x))$, and on
  rejection the replacement is drawn from the residual
  $p'(x) = \max(0, p(x)-q(x)) \,/\, \sum_{x'} \max(0, p(x')-q(x'))$.

  This is the nicest result in the lecture, so slow down for it.
  It is Metropolis--Hastings-style rejection sampling applied
  to token distributions, and the residual term is exactly what repairs the
  bias introduced by proposing from the wrong distribution.

  The practical consequence is what makes it unusual among everything else
  today: every other technique in this lecture trades quality for speed, and
  you have to measure what you broke. This one cannot break anything. The only
  risk is that it fails to help.

  Caveat for the pedants, and worth saying because it returns in section 8:
  ``exact'' means exact in exact arithmetic. In floating point, $p$ computed over
  a batch of 5 positions is not bitwise identical to $p$ computed over 1, so
  speculative decoding does change your outputs at the bit level. Hold that
  thought.}
\end{frame}

% Optional quantitative speculative-decoding detail.
\iffalse
\begin{frame}[shrink=16]{Acceptance rate decides whether this pays}
  The \key{acceptance rate} is the fraction of the drafter's guesses the big
  model keeps, and it decides whether any of this is worth doing.
  With 4 guesses per round:

  \vspace{0.3em}
  \centering
  \begin{tabular}{@{}C{3.0cm} C{3.2cm} C{3.0cm}@{}}
    \toprule
    \textbf{acceptance rate} & \textbf{tokens per big-model pass} &
      \textbf{net speed-up}\\
    \midrule
    0.8 & 3.36 & $\approx 2.4\times$\\
    0.6 & 2.31 & $\approx 1.6\times$\\
    0.5 & 1.94 & $\approx 1.4\times$\\
    0.3 & 1.42 & $\approx 1.0\times$ --- no gain\\
    \bottomrule
  \end{tabular}

  \vspace{0.5em}
  \raggedright\small
  Net speed-up assumes a drafter costing about 10\% of a big-model pass, so
  4 guesses add about 0.4 passes of overhead.
  \vspace{0.6em}

  \begin{block}{What moves the acceptance rate}
    \small
    It goes up with a drafter trained on the same data with the same tokenizer,
    a domain the drafter knows, and a \key{low temperature}, since easy tokens
    are easier to agree on. Unfamiliar text, a high temperature or a tokenizer
    mismatch will collapse it. \key{Measure it on your own data.}
  \end{block}
  \source{Leviathan et al., ICML 2023.}
  \note{The closed form behind the table: with per-token acceptance
  probability $\alpha$ and $k$ proposals, the expected number of tokens per
  target pass is $(1 - \alpha^{k+1})/(1-\alpha)$ --- a truncated geometric
  series, because you keep accepting until the first rejection. Every row uses
  $k=4$.

  The table is the point: below about 0.4 acceptance the whole scheme
  stops paying for itself, and you have added a second model, more memory and
  more complexity for nothing.

  Note also the interaction with batching, which is the part the papers
  underplay. Speculative decoding spends surplus compute, and under heavy
  continuous batching you no longer \emph{have} any, because the arithmetic
  units are busy serving other users. So it helps most at low concurrency and
  can hurt at high concurrency, and good servers turn it off adaptively.

  So it is a latency optimisation, which is the opposite of what continuous
  batching gives you.}
\end{frame}

% Optional current systems detail.
\begin{frame}[shrink=3]{Self-speculation: no second model}
  \begin{columns}[T]
    \begin{column}{0.5\textwidth}
      Instead of a second small model, let the big model guess ahead for
      itself.
      \vspace{0.4em}

      \textbf{Medusa} --- bolt a few extra output heads onto the model so it
      guesses the next two or three tokens at once, then verify them all.
      \vspace{0.4em}

      \textbf{EAGLE} --- guess using the model's own internal states rather
      than its output tokens. Higher acceptance for the same overhead; the
      current default in several stacks.
      \vspace{0.4em}

      \textbf{MTP heads} --- ``multi-token prediction'' heads trained
      \emph{during pretraining} as a side objective, then reused as a free
      drafter at inference.
    \end{column}
    \begin{column}{0.5\textwidth}
      \begin{block}{MTP in 2026 models}
        \small
        \textbf{DeepSeek V3} --- MTP-1 during training, optional at inference.\\[0.3em]
        \textbf{GLM-4.7, MiniMax M2.1} --- MTP-3.\\[0.3em]
        \textbf{Step 3.5 Flash} --- MTP-3 at training \emph{and} inference,
        which is unusual; MTP is normally a training-only trick.\\[0.3em]
        \textbf{Nemotron 3 Super} --- shared-weight MTP, explicitly for
        speculative decoding.\\[0.3em]
        \textbf{Kimi K2.5, Arcee Trinity} --- no MTP at all.
      \end{block}
      \vspace{0.2em}
      \small
      Compared with a separate drafter you get the same tokenizer, the same
      training data and almost no extra memory.
    \end{column}
  \end{columns}
  \sources{Cai et al., \textit{Medusa}, 2024; Li et al., \textit{EAGLE}, 2024;
    Raschka, \textit{A Visual Guide to Attention Variants}, Mar 2026 and
    \textit{A Dream of Spring for Open-Weight LLMs}, Feb 2026, for the
    adoption map.}
  \note{Mechanics, if asked: Medusa attaches extra decoding heads to the final
  hidden state, predicting positions $t{+}1, t{+}2, \dots$ in parallel and
  verifying a tree of candidates; EAGLE drafts in \emph{feature} space using
  the target's own hidden states, which is why its acceptance rate is higher.

  The through-line: MTP started as a way to get a richer training signal
  --- predict several tokens ahead, sum the losses --- and was described as
  training-only. Then the inference people noticed that a training-time
  auxiliary head is a free drafter with a perfectly matched tokenizer. Nice
  example of a technique acquiring a second life.

  Point at the ``no MTP'' row too. Two 2026 flagships declined it, which tells
  you the gain is real but not overwhelming, and it costs training complexity.

  Practical note: for their 4B project models this is mostly out of reach in a
  90-minute lab, so it is context rather than coursework. Say so.}
\end{frame}
\fi

\begin{frame}{Constrained decoding}
  \begin{columns}[T]
    \begin{column}{0.5\textwidth}
      \textbf{What it is}
      \begin{itemize}
        \item Give the server a schema or grammar for the required output.
        \item At each step, mask tokens that would make the output invalid.
        \item Sample from the tokens that remain.
      \end{itemize}
      \vspace{0.3em}
      \textbf{Why it is used}\\
      Applications often need parseable JSON, and retries add latency and
      cost.
    \end{column}
    \begin{column}{0.5\textwidth}
      \begin{block}{What it guarantees}
        \small
        \textbf{Valid syntax when generation completes successfully.}
        The constraint removes invalid token sequences, but truncation and
        server failures still need handling.
      \end{block}
      \begin{block}{What it does not guarantee}
        \small
        \textbf{Anything about the content.} The label inside the valid JSON
        can still be wrong, and the constraint can make it wronger.
      \end{block}
    \end{column}
  \end{columns}
  \vspace{0.4em}
  Implementations such as \textbf{Outlines} and \textbf{XGrammar} make the
  allowed-token lookup efficient.
  \sources{Willard \& Louf, \textit{Efficient Guided Generation for LLMs},
    2023 (Outlines); Dong et al., \textit{XGrammar}, 2024.}
  \note{The FSM compilation is the clever part: because the automaton is over
  the \emph{tokenizer's} vocabulary rather than over characters, the allowed
  set can be precomputed per state, and the run-time cost drops to a mask
  lookup.

  Note how closely this resembles min-$p$ and top-$k$ from section 1: it is the
  same operation, masking logits before the softmax. The only difference is
  that the mask comes from a grammar rather than from the distribution itself.

  Set up the gotcha explicitly by reading the two blocks as a pair, then turn
  the page.}
\end{frame}

\begin{frame}{Valid syntax can still reduce task accuracy}
  A grammar guarantees the \emph{shape} of an output, not its content.
  The model may prefer to reason first, use another key order or decline to
  answer. Masking those continuations can force it onto a low-probability path
  and reduce answer quality.
  \vspace{0.4em}
  \centering
  \begin{tabular}{@{}L{4.0cm} C{3.0cm} C{3.0cm}@{}}
    \toprule
    & \textbf{invalid output} & \textbf{task score}\\
    \midrule
    free generation + parser & 4--15\% & baseline\\
    JSON-schema constrained & \good{0\%} & \warn{report it and see}\\
    \bottomrule
  \end{tabular}

  \vspace{0.3em}
  \raggedright
  Report \textbf{both} invalid-output rate and task score. For reasoning
  models, let the model reason freely and constrain only the final answer.
  \note{Gotcha two of three, and it is Thursday's ``constrain'' step in a single
  slide. The lab deliberately asks for both numbers so that the trade-off is
  visible rather than assumed.

  The mechanism to explain is distribution shift: the model is being forced
  onto conditional paths it has low probability for, and a low-probability
  prefix tends to be followed by low-quality continuations. A schema that
  demands the answer field first is actively preventing chain-of-thought.

  The mitigation in the last line is what a good engineer does and it costs
  nothing: unconstrained reasoning, constrained answer. Tell them to try both
  orders on Thursday if they have time.

  Clock check: should be around 17:20 at the end of this section.}
\end{frame}

% =====================================================================
\section{Serving systems and cost}
% =====================================================================

\begin{frame}{When one GPU is not enough}
  First identify the constraint that one GPU cannot meet.
  \vspace{0.3em}
  \centering
  \begin{tabular}{@{}L{4.2cm} L{3.0cm} L{4.0cm}@{}}
    \toprule
    \textbf{The problem} & \textbf{The answer} & \textbf{What it costs}\\
    \midrule
    weights and KV do not fit, or decode is too slow &
      \key{split the model} & communication on every token\\[0.3em]
    they fit, throughput is too low &
      \key{more replicas} & more hardware, little communication\\[0.3em]
    MoE experts dominate memory &
      expert parallel & every GPU talks to every other\\
    \bottomrule
  \end{tabular}

  \vspace{0.4em}
  \raggedright
  If the model fits and one copy is fast enough, run several independent
  replicas before splitting one request across GPUs. Replicas do not exchange
  activations and are easier to scale and recover.
  \note{Clock check: should be around 17:20--17:22.

  The names for the cost column, which the next three slides use: tensor
  parallelism pays an \emph{all-reduce} (every GPU's partial result summed and
  broadcast back) twice per transformer block; pipeline parallelism pays a
  \emph{bubble} (idle stages); expert parallelism pays an \emph{all-to-all}
  per MoE layer.

  Lead with the replica rule because it is the one that applies to everyone in
  this room: a 4B model on a 96\,GiB card never needs tensor parallelism, and
  someone will try it anyway and get a slower server.

  The three-row structure is the decision tree. The useful question is which
  constraint you are hitting: memory, latency or throughput. Each has a
  different answer, and people confuse them all the time.}
\end{frame}

% Detailed multi-GPU mechanisms are optional reference material. The
% projected lecture keeps the decision slide that explains when they are used.
\iffalse
\begin{frame}{Tensor parallelism}
  \begin{columns}[T,onlytextwidth]
    \begin{column}{0.47\textwidth}
      \centering
      \begin{tikzpicture}[node distance=4mm, font=\scriptsize]
        \node[modmute] (x) {activations $x$};
        \node[mod, below left=6mm and 0mm of x] (g0) {GPU 0\\\tiny half the heads,\\half the FFN};
        \node[mod, below right=6mm and 0mm of x] (g1) {GPU 1\\\tiny half the heads,\\half the FFN};
        \node[modhi, below=14mm of x] (ar) {all-reduce};
        \node[modmute, below=4mm of ar] (y) {output};
        \draw[flow] (x) -- (g0);
        \draw[flow] (x) -- (g1);
        \draw[flowhi] (g0) -- (ar);
        \draw[flowhi] (g1) -- (ar);
        \draw[flow] (ar) -- (y);
      \end{tikzpicture}
    \end{column}
    \begin{column}{0.53\textwidth}
      Every weight matrix is split across the GPUs, so each one does part of
      every layer. Attention splits by head and the FFN by hidden dimension.
      \begin{itemize}\small
        \item \key{two all-reduces per transformer block}, on every token:
          each GPU's partial result is summed and sent back to all of them
        \item the KV cache is split too, so concurrency scales as well
      \end{itemize}
      \begin{block}{The condition}
        \small
        You need a \key{fast interconnect}, meaning NVLink inside a node. Over
        Ethernet it is usually a net loss, because the all-reduce sits on the
        critical path of every decode step.
      \end{block}
    \end{column}
  \end{columns}
  \source{Shoeybi et al., \textit{Megatron-LM}, 2019;
    Pope et al., \textit{Efficiently Scaling Transformer Inference}, MLSys 2023.}
  \note{Contrast with week 2 deliberately: in training, communication can hide
  behind a large batch's compute. In decode there is almost no compute to hide
  behind, so the all-reduce is naked latency, twice per layer, on every token.
  That is why interconnect quality decides whether tensor parallelism helps at
  all.

  Say the upside out loud, because it is no longer on the slide: tensor
  parallelism is the only option here that reduces per-token \emph{latency},
  because it divides the bytes each GPU has to read. Two GPUs, half the weights
  each, roughly half the TPOT.

  Roihu nodes have four GH200s with a fast interconnect, so a student who
  genuinely needed it could use TP=2. Nobody on this course will.}
\end{frame}

\begin{frame}{Pipeline parallelism and the bubble}
  \begin{columns}[T,onlytextwidth]
    \begin{column}{0.5\textwidth}
      Split the model \emph{by layer}: GPU 0 holds layers 1--18, GPU 1 holds
      19--36. Activations are passed forward once per stage.
      \begin{itemize}\small
        \item communication is tiny: one activation tensor rather than an
          all-reduce
        \item works over commodity interconnect
        \item memory scales with the number of stages
      \end{itemize}
      \vspace{0.4em}
      \key{But:} with one request in flight, stage 1 is idle while stage 0
      works. That idleness is the \textbf{bubble}.
    \end{column}
    \begin{column}{0.5\textwidth}
      \centering
      \begin{tikzpicture}[x=0.82mm, y=1mm, font=\tiny]
        \draw[UHGrey, ->] (0,-4) -- (62,-4) node[lbl, anchor=west] {time};
        \node[lbl, anchor=east] at (-2,13) {stage 0};
        \node[lbl, anchor=east] at (-2,4) {stage 1};
        \foreach \i in {0,1,2,3} {
          \fill[UHBlue!35, draw=UHBlue] ({\i*14},10) rectangle ({\i*14+12},16);
          \fill[UHOrange!25, draw=UHOrange] ({\i*14+14},1) rectangle ({\i*14+26},7);
        }
        \fill[UHGrey!20, draw=UHGrey!40] (0,1) rectangle (12,7);
        \node[text=UHGrey] at (6,4) {bubble};
      \end{tikzpicture}

      \vspace{0.6em}
      \raggedright\small
      The fix is the same as in training: keep many micro-batches in flight.
      In serving, continuous batching supplies them for free, as long as there
      is enough traffic.
    \end{column}
  \end{columns}
  \vspace{0.5em}
  \centering
  \small
  Rule of thumb: \key{tensor parallel inside a node, pipeline parallel between
  nodes.}
  \note{The comparison to make is communication volume against idle time.
  Tensor parallelism communicates a lot but keeps every GPU busy, while
  pipeline parallelism barely communicates and then wastes time whenever load
  is low. Which you prefer depends on your interconnect and your traffic.

  The serving-specific point is the nice one: in training you have to
  manufacture micro-batches to fill the pipeline, whereas a busy inference
  server already has a queue of independent requests. Pipeline parallelism
  therefore behaves much better under load than at low load, which is the
  opposite of the intuition most people bring.

  The rule of thumb at the bottom is what production systems actually do, and
  it follows directly from where the fast links are.}
\end{frame}

\begin{frame}{Expert parallelism, for MoE}
  \begin{columns}[T]
    \begin{column}{0.54\textwidth}
      An MoE model's parameters are mostly experts, and at inference only
      top-$k$ of them run per token. Putting \emph{different experts on
      different GPUs} is the natural split.
      \vspace{0.5em}

      The communication pattern is an \key{all-to-all}: every GPU sends each
      token to whichever GPUs hold its chosen experts, then collects the
      results back.
      \vspace{0.5em}

      \textbf{Load balance turns into a serving problem too.} If a batch's
      routing happens to favour a few experts, the GPUs holding them become
      the critical path and the rest wait.
    \end{column}
    \begin{column}{0.46\textwidth}
      \begin{block}{Why total parameters still matter}
        \small
        Qwen3-Next 80B-A3B activates 3B parameters per token, but you still
        have to \emph{hold} all 80B. At bf16 that is 160\,GB, so it does not
        fit on one GH200 at any batch size.
        \vspace{0.3em}
        \key{You get 3B speed on 80B worth of hardware.} That is week 1's
        gotcha with a price attached.
      \end{block}
      \vspace{0.2em}
      \small
      It is also why some 2026 teams deliberately chose \emph{coarser}
      experts: fewer and larger experts mean less routing overhead and simpler
      placement.
    \end{column}
  \end{columns}
  \source{Raschka, \textit{A Dream of Spring for Open-Weight LLMs}, Feb 2026,
    on MoE granularity; DeepSeek V3 technical report on expert parallel
    deployment.}
  \note{Reconnect to the week 1 MoE gotcha explicitly. Counting active
  parameters tells you about speed and nothing at all about memory, and here
  is what that costs in euros: an 80B-A3B model runs at the speed of a 3B
  model while taking up the room of an 80B one, and you pay for both of those
  properties separately.

  The load-balance point surprises people who have only read the training
  literature. Auxiliary load-balancing losses shape the \emph{average} over a
  run, but a single live batch can still be badly skewed, and then the busiest
  expert sets your tail latency.

  Keep this brief, since nobody in the room will serve an MoE this term. It is
  context for reading deployment reports.}
\end{frame}

\begin{frame}{Disaggregated prefill and decode}
  \centering
  \begin{tikzpicture}[x=0.9mm, y=1mm, font=\scriptsize]
    \node[modmute] (req) at (-4,10) {requests};
    % prefill pool
    \draw[draw=UHBlue, fill=UHBlue!8, rounded corners=2pt] (18,0) rectangle (52,22);
    \node[lbl, anchor=south] at (35,22.5) {prefill pool: limited by arithmetic};
    \node[mod, minimum width=12mm] at (27,15) {GPU};
    \node[mod, minimum width=12mm] at (43,15) {GPU};
    \node[lbl] at (35,5) {big matmuls, FP8, high MFU};
    % decode pool
    \draw[draw=UHOrange, fill=UHOrange!8, rounded corners=2pt] (74,0) rectangle (116,22);
    \node[lbl, anchor=south] at (95,22.5) {decode pool: limited by bandwidth};
    \node[mod, minimum width=11mm] at (82,15) {GPU};
    \node[mod, minimum width=11mm] at (95,15) {GPU};
    \node[mod, minimum width=11mm] at (108,15) {GPU};
    \node[lbl] at (95,5) {huge KV pool, large batches};
    % arrows
    \draw[flow] (req) -- (17,10);
    \draw[flowhi] (52,11) -- (73,11);
    \node[lbl, anchor=south, text=UHOrange] at (62.5,12) {KV transfer};
    \node[lbl, anchor=north] at (62.5,9) {NVLink / RDMA};
    \draw[flow] (116,11) -- (128,11);
    \node[lbl, anchor=west] at (128,11) {tokens};
  \end{tikzpicture}

  \vspace{1.0em}
  \raggedright
  \begin{columns}[T]
    \begin{column}{0.54\textwidth}
      \small
      Two phases with opposite bottlenecks were sharing one pool and
      interfering with each other. Separate them and you can:
      \begin{itemize}
        \item size each pool independently for its actual load
        \item tune parallelism and even \emph{hardware} per phase
        \item stop long prefills from touching anyone's TPOT at all
      \end{itemize}
    \end{column}
    \begin{column}{0.46\textwidth}
      \begin{block}{The cost}
        \small
        The KV cache for the whole prompt must cross the wire before decoding
        starts. That is tens to hundreds of megabytes per request, which is why
        this needs NVLink or RDMA and why it only pays at scale.
      \end{block}
    \end{column}
    \end{columns}
  \source{Zhong et al., \textit{DistServe}, OSDI 2024; now supported in vLLM
    and SGLang.}
  \note{This is how large deployments are built now, and it follows from
  section 2. Prefill saturates the arithmetic units while decode spends its
  time waiting on memory, so putting both on the same machines guarantees that
  one of them is wasting whatever the other needs.

  The forward-looking consequence is more interesting. Once the phases are
  separated there is no reason they need the same \emph{kind} of GPU: a
  prefill box wants FLOPs more than memory, and a decode box wants the
  opposite. Expect the hardware to diverge.

  Say clearly that this is well beyond anything they will do on Roihu. It is
  here so they can read a serving architecture diagram and know what they are
  looking at.}
\end{frame}
\fi

% Engine catalogues change quickly; the lab uses vLLM and the reading version
% retains the comparison.
\iffalse
\begin{frame}[shrink=19]{Four serving engines and what each is for}
  \centering
  \footnotesize
  \begin{tabular}{@{}L{2.2cm} L{4.4cm} L{4.6cm}@{}}
    \toprule
    \textbf{Engine} & \textbf{Good at} & \textbf{Choose it when}\\
    \midrule
    \textbf{vLLM} & PagedAttention, continuous batching, huge model coverage,
      OpenAI-compatible API, LoRA adapters served live &
      default for research and for most production; \textbf{this course}\\[0.3em]
    \textbf{SGLang} & RadixAttention prefix reuse, structured and multi-call
      programs, deterministic-inference mode &
      agent loops, heavy prefix sharing, reproducibility\\[0.3em]
    \textbf{TensorRT-LLM} & compiled kernels, the best raw numbers on NVIDIA
      hardware &
      a fixed model you will serve for months and can afford to compile\\[0.3em]
    \textbf{llama.cpp / Ollama} & CPU and consumer GPUs, GGUF, aggressive
      low-bit quantisation, single-user &
      a laptop, an edge device, a demo\\
    \bottomrule
  \end{tabular}

  \vspace{0.8em}
  \raggedright
  \begin{block}{Caveat worth carrying out of this course}
    \small
    Engine support is not uniform across architectures. Classic GQA models have
    years of tuned kernels behind them, and the hybrid, linear and sparse
    attention models usually have very few. \key{In late 2026, a model that is
    cheaper on paper can still be slower in your stack.} Benchmark the model
    and the engine together rather than the model alone.
  \end{block}
  \note{The honest summary of the field is that vLLM and SGLang keep absorbing
  each other's features, so the distinction narrows with every release. Pick on
  ecosystem fit rather than on a benchmark from six months ago.

  llama.cpp is on the slide for a reason that has nothing to do with
  performance. It is how most people outside this room actually run a local
  model, and GGUF plus Ollama is what their non-specialist colleagues will ask
  them about.

  The caveat block is the practitioner's point from the course outline, and it
  generalises. A paper FLOP count is a prediction and the kernels are what you
  actually get, and most disappointment in this area lives in that gap. It is
  also why I would still default to a boring GQA model for a deployment I had
  to support.}
\end{frame}

\fi
\begin{frame}{Cost per million tokens, from first principles}
  \small
  \[
    \text{EUR per million tokens} \;=\;
    \frac{\text{EUR per GPU-hour}}{\text{tokens you got out of that hour}}
    \times 10^{6}
  \]
  \emph{In words:} you rent the card by the hour, so divide the rent by what
  the hour produced. Rent is EUR 2.50/h and the card is busy only 40\% of the
  time (its \key{utilisation}), so everything divides by 0.4 at the end.
  \begin{columns}[T]
    \begin{column}{0.5\textwidth}
      \textbf{32 users at once, measured 2{,}500 tokens/s:}
      \begin{itemize}
        \item $2500 \times 3600 = 9$ million tokens an hour
        \item EUR 2.50 for 9 million $=$ EUR 0.28 per million
        \item busy only 40\% of the time: \key{EUR 0.69 per million}
      \end{itemize}
      \textbf{One user, measured 120 tokens/s:} 0.43 million tokens an hour,
      the same arithmetic, \key{EUR 14.5 per million}.
    \end{column}
    \begin{column}{0.5\textwidth}
      \begin{block}{The same GPU, the same model, $21\times$ the price}
        \small
        Nothing changed except how many requests were in flight. Batch size is
        the dominant term in your unit economics.
      \end{block}
      \vspace{0.3em}

      Public prices in late 2026: roughly \$0.02--0.10 per million output
      tokens for hosted 4B-class open models, \$5--15 for frontier ones.
    \end{column}
  \end{columns}
  \source{GPU price: representative on-demand H100/GH200 cloud rate, late 2026.
    Throughput: the kind of figure Thursday's lab produces. API prices: public
    price lists, November 2026 --- \textbf{these rot fast}.}
  \note{Do this arithmetic on the board. It is four numbers and a division, and
  it is probably the most immediately employable thing in the lecture. The
  general form, with $u$ for utilisation, is
  EUR per $10^{6}$ tokens $=$ (EUR per GPU-hour) / (tokens/s $\times$ 3600
  $\times\ u$) $\times\ 10^{6}$; the single-user case is
  $2.50 / 4.32\times10^{5} \times 10^{6} \div 0.4 \approx 14.5$.

  The $21\times$ is the punchline and it reframes everything before it:
  sections 3 and 4 were all about moving from the second calculation to the
  first.

  Two framings that used to be on the slide and are worth saying anyway. The
  gap between your number and a provider's is utilisation and batch size
  rather than some secret they have, and batching is not a performance detail
  at the margin: it is the largest term in the whole calculation.

  Be precise about the utilisation factor, because people leave it out and then
  wonder why their estimate is optimistic. You pay for the GPU every hour it is
  allocated, including the nights and the idle afternoons. A 40\% duty cycle is
  realistic for a service with human traffic; a batch job can approach 100\%
  and is correspondingly cheap.

  Date the API prices out loud and say they will be wrong by spring.}
\end{frame}

\begin{frame}{Buying hardware for a decode-heavy workload}
  \centering
  \begin{tabular}{@{}L{3.6cm} R{2.4cm} R{2.4cm} R{2.4cm}@{}}
    \toprule
    & \textbf{bf16 TFLOP/s} & \textbf{bandwidth} & \textbf{memory}\\
    \midrule
    GH200 (96\,GiB, HBM3)  & $\approx$ 990 & 4.0\,TB/s & 96\,GiB\\
    GH200 (144\,GB, HBM3e) & $\approx$ 990 & 4.9\,TB/s & 144\,GB\\
    A100 80\,GB            & $\approx$ 312 & 2.0\,TB/s & 80\,GB\\
    L40S                   & $\approx$ 362 & 0.86\,TB/s & 48\,GB\\
    \bottomrule
  \end{tabular}

  \vspace{0.4em}
  \raggedright
  \small
  Compute and bandwidth are the two halves of the roofline. The L40S has more
  advertised arithmetic than an A100 but less than half its memory bandwidth.
  A decode-heavy workload can therefore run faster on the A100.
  \vfill
  \begin{takeaway}
    For a decode-heavy workload, compare \key{bandwidth per euro} and memory
    capacity, not only peak arithmetic.
  \end{takeaway}
  \note{Clock check: should be around 17:34.

  The A100-versus-L40S comparison is the whole slide. It is a concrete case
  where the number everyone quotes gives you the wrong answer. Ask the room
  which they would buy before showing the last line.

  Tie it back to the first-principles formula from section 2: tokens per second
  is bandwidth divided by bytes per token, so for a single stream bandwidth
  is a useful first predictor of throughput. Everything else is a correction
  term.}
\end{frame}

% =====================================================================
\section{Research hook: deterministic inference}
% =====================================================================

% This case study is useful reading but too specialised for the core
% 90-minute inference overview.
\iffalse
\begin{frame}[shrink=18]{The demonstration}
  Set \texttt{temperature=0}, \texttt{seed=0}, the same prompt, the same
  server, the same weights.
  \vspace{0.6em}
  \centering
  \begin{tabular}{@{}L{6.2cm} L{5.0cm}@{}}
    \toprule
    \textbf{Experiment} & \textbf{Result}\\
    \midrule
    Send the prompt twice, alone, back to back &
      \good{bitwise identical}\\[0.3em]
    Send it alone, then send it inside a batch of 64 &
      \bad{different output}\\[0.3em]
    Send it 1{,}000 times under live traffic &
      \bad{many distinct completions}\\
    \bottomrule
  \end{tabular}

  \vspace{0.8em}
  \raggedright
  A published run of this experiment: 1{,}000 completions of 1{,}000 tokens
  each from one model at temperature 0 produced \key{80 distinct completions},
  with the first divergence typically around the \key{103rd token}.
  \vspace{0.6em}
  \begin{gotcha}[Temperature 0 is not deterministic]
    And the cause is \textbf{not} GPU randomness, \textbf{not} a missing seed,
    and \textbf{not} your sampler. It is \key{who else was in your batch}.
  \end{gotcha}
  \source{Thinking Machines Lab, \textit{Defeating Nondeterminism in LLM
    Inference}, September 2025.}
  \note{Clock check: should be around 17:34--17:35 here. Ten minutes left, and
  this section is worth every one of them.

  Run the demonstration as a story, not as a slide. Everyone in the room
  believes temperature 0 means argmax means deterministic, because that is what
  the maths says. Row one confirms their belief. Row two destroys it.

  Then kill the three wrong explanations immediately and by name, because
  otherwise people spend the next five minutes privately theorising about
  \texttt{atomicAdd} and cosmic rays. The seed is irrelevant: at temperature 0
  nothing is sampled. GPU atomics are a real source of run-to-run
  nondeterminism in \emph{training}, but they are not what is happening here,
  and a forward pass run twice on identical inputs really is reproducible.

  The number, 80 distinct completions out of 1{,}000, is public and citable.
  Use it, because it makes the problem feel like engineering rather than
  folklore.}
\end{frame}

\begin{frame}[shrink=5]{Floating-point addition is not associative}
  In IEEE double precision:
  \[
    (0.1 + 0.2) + 0.3 = 0.6000000000000001
    \quad\neq\quad
    0.1 + (0.2 + 0.3) = 0.6
  \]
  \vspace{-0.4em}
  \begin{columns}[T]
    \begin{column}{0.54\textwidth}
      Each addition rounds to the nearest representable value, so the rounding
      error depends on \key{the order in which you add}. Add the same numbers
      in a different order and you get a slightly different sum.
      \vspace{0.5em}

      And we are not even in double precision. A bf16 number keeps only
      \key{8 bits of precision}, so each rounding step is about 0.4\% of the
      value, against roughly one part in $10^{16}$ for a double. It is the
      same effect, about thirteen orders of magnitude larger.
    \end{column}
    \begin{column}{0.46\textwidth}
      \begin{block}{From one bit to a different answer}
        \small
        A change in summation order perturbs a token's score in its last bits.
        \vspace{0.3em}
        That changes \emph{which} token is top only when two candidates are
        nearly tied, which is rare per token, \textbf{but you generate
        thousands of tokens}.
        \vspace{0.3em}
        One flip, and every token after it is conditioned on a different
        prefix. \key{The two generations fork rather than drifting apart.}
      \end{block}
    \end{column}
  \end{columns}
  \vspace{0.3em}
  \centering
  \small
  Nothing here is random. Every run is deterministic \emph{given the order in
  which the additions happened}.
  \note{Precisely: bf16 keeps 8 significand bits, so its relative rounding step
  is about $2^{-8} \approx 0.4\%$, against $2^{-53}$ for an IEEE double.
  ``Which token is top'' is the $\arg\max$ over the logits.

  The 0.1/0.2/0.3 example is worth typing into a Python prompt live if the
  projector allows. It takes ten seconds and it is far more convincing than a
  slide.

  The conceptual move to land is in the last line. None of this is
  nondeterminism in the probabilistic sense. Every kernel is a deterministic
  function of its inputs \emph{and of how it was scheduled}. What changes is
  the scheduling, and that is the thing we are not controlling.

  The fork metaphor matters because it explains why the divergence is so
  visible. Models do not drift apart slowly. One token differs, and from then
  on the two generations are conditioned on different text, so they can end up
  saying opposite things.}
\end{frame}

\begin{frame}[shrink=18]{Where the summation order comes from}
  \centering
  \begin{tikzpicture}[x=1mm, y=1mm, font=\scriptsize]
    \node[lbl, anchor=east] at (-2,22) {batch of 1};
    \foreach \i in {0,...,7}
      \node[tokcell, minimum width=6mm] at ({6+\i*8},22) {};
    \draw[decorate, decoration={brace, amplitude=3pt}, UHBlue, thick]
      (3,27) -- (69,27);
    \node[lbl, anchor=south, text=UHBlue] at (36,28) {8 numbers to add, done in 2 pieces};

    \node[lbl, anchor=east] at (-2,6) {batch of 64};
    \foreach \i in {0,...,7}
      \node[tokcell, minimum width=6mm, fill=UHOrange!15] at ({6+\i*8},6) {};
    \draw[decorate, decoration={brace, amplitude=3pt}, UHOrange, thick]
      (3,11) -- (69,11);
    \node[lbl, anchor=south, text=UHOrange] at (36,12)
      {the same 8 numbers, now in 4 pieces, so different partial sums};
  \end{tikzpicture}

  \vspace{1.2em}
  \raggedright
  To add up a long list of numbers fast, a GPU chops it into pieces, adds each
  piece separately, then adds the pieces. \key{How many pieces} is chosen at
  run time from the shape of the work, and that shape includes everyone else's
  requests:
  \begin{itemize}\small
    \item \key{matrix multiply}: the chunking is picked from the batch size, so
      a different batch size gives a different chunking and a different order
      of additions
    \item \key{attention}: the standard fast kernel splits the cache into a
      \emph{fixed number} of pieces, computed from the batch and the sequence
      length, so both change the order
    \item \key{chunked prefill and prefix caching} change how much of the
      sequence is processed in one pass, so a \textbf{cache hit can change
      your output}
  \end{itemize}
  \key{Your token scores depend on other users' requests.}
  \source{Thinking Machines Lab, Sept 2025; FlashAttention / FlashDecoding
    split-KV scheduling, Dao et al. 2022--2023.}
  \note{This is the technical core of the section. Go slowly. The proper names
  for the three bullets: matmul tile sizes and split-K factors chosen from the
  batch dimension; FlashDecoding's split-KV, which divides the reduction over
  the KV length into a fixed number of chunks; and scheduler-level effects
  from chunked prefill and prefix caching.

  The split-K and split-KV detail is the one that makes it click. These kernels
  choose \emph{how many} parallel chunks to split a reduction into so as to
  keep all the streaming multiprocessors busy. With one sequence in flight
  there is little work, so they split aggressively to fill the GPU. With
  sixty-four there is plenty of work, so they split less. Either choice is
  correct and fast, and they produce different last bits.

  The third bullet connects back to today's own lecture and usually gets a
  reaction: prefix caching, which I recommended two sections ago as a free win
  with no quality cost, changes your output bits. Enjoy that tension. It is
  honest, and it is the sort of interaction that only shows up once you have
  built the system.}
\end{frame}

\begin{frame}[shrink=19]{A better name for the problem: batch invariance}
  \begin{block}{Definition}
    A kernel is \key{batch-invariant} if the result computed for a given
    request is \textbf{bitwise identical} regardless of the size and
    composition of the batch it was executed in, and regardless of how the
    work was split.
  \end{block}
  \vspace{0.5em}
  \begin{columns}[T]
    \begin{column}{0.5\textwidth}
      \textbf{Why this framing is better.}
      \begin{itemize}\small
        \item ``GPU nondeterminism'' suggests randomness. There is none.
        \item ``Set a seed'' fixes nothing, because nothing is being sampled.
        \item ``Run-to-run reproducibility'' aims at the wrong thing. The
          server already has it \emph{at fixed batch composition}.
      \end{itemize}
      \vspace{0.4em}
      What you want is invariance to a variable nobody ever meant to be an
      input.
    \end{column}
    \begin{column}{0.5\textwidth}
      \begin{block}{Three kernels have to change}
        \small
        \textbf{RMSNorm}: fix the reduction strategy instead of choosing it
        from the batch size.\\[0.3em]
        \textbf{Matmul}: fix the tile and split-K configuration so the kernel
        selector cannot pick one per shape.\\[0.3em]
        \textbf{Attention} is the hard one. Split the KV reduction into chunks
        of \key{fixed size} rather than a fixed \emph{number}, so the
        summation tree does not depend on sequence length or batch.
      \end{block}
    \end{column}
  \end{columns}
  \source{Thinking Machines Lab, \textit{Defeating Nondeterminism in LLM
    Inference}, Sept 2025.}
  \note{Naming the property correctly is most of the contribution here, and
  there is a transferable research lesson in that. If you frame the problem as
  ``the GPU is flaky'' there is nothing to do about it. Frame it as ``this
  function has an argument nobody declared'' and the fix is obvious.

  The fixed-size versus fixed-count distinction in the attention bullet is the
  key engineering idea. A fixed \emph{number} of splits adapts to load and runs
  faster. A fixed chunk \emph{size} means a given sequence is always summed the
  same way, whatever else is on the GPU, and you give up some speed for that.

  If anyone asks why this was not built in from the start: nobody was asking
  for it, and the adaptive version is genuinely faster. It was a reasonable
  default that became a problem once people started doing RL with these
  servers.}
\end{frame}

\begin{frame}[fragile,shrink=8]{Making it real, and what it costs}
\begin{lstlisting}[language=slurm]
python -m sglang.launch_server \
  --model-path Qwen/Qwen3.5-4B-Instruct \
  --enable-deterministic-inference \
  --attention-backend flashinfer
\end{lstlisting}
  \vspace{0.4em}
  \begin{columns}[T]
    \begin{column}{0.54\textwidth}
      \small
      What a deterministic mode has to do:
      \begin{itemize}
        \item pin batch-invariant kernels for RMSNorm, the linear layers and
          attention
        \item make the attention backend's KV splitting independent of batch
          and sequence length
        \item keep the \emph{scheduler} from changing what a request sees, so
          chunked prefill boundaries and prefix-cache hits do not alter the
          computation
      \end{itemize}
      \vspace{0.3em}
      That last point is what makes this a serving-systems problem rather than
      a kernel problem.
    \end{column}
    \begin{column}{0.46\textwidth}
      \begin{block}{The price}
        \small
        Published baseline for the first implementation: a benchmark that took
        \textbf{26\,s} unmodified took \textbf{55\,s} with naive deterministic
        kernels, improving to \textbf{42\,s} once attention was done properly.
        That is roughly a \key{$1.6\times$} cost.
        \vspace{0.3em}
        Engine work since then has been about closing that gap. The honest
        answer today is tens of percent, depending on backend and batch size,
        and nobody should tell you it is free.
      \end{block}
    \end{column}
  \end{columns}
  \sources{Thinking Machines Lab, Sept 2025, for the published timings;
    SGLang deterministic-inference support. Figures are backend- and
    workload-dependent, so measure your own.}
  \note{This is my own area, so be concrete and be careful in equal measure.

  The honest status: the kernel side is now well understood and mostly a
  matter of discipline. The interesting remaining work is on the scheduler
  side, making the guarantee survive chunked prefill, prefix caching,
  speculative decoding and multi-GPU reductions. All of those exist to go fast
  by adapting to load, which is exactly what we are now forbidding.

  Be explicit about the cost and do not round it away. Anyone who tells you
  determinism is free has not measured it under load. For practitioners the
  useful framing is that it is a \emph{mode} you turn on for evaluation,
  debugging and RL rollouts, and leave off when you are serving users who do
  not care.

  Do not overclaim my role. Say ``this is the line of work I am involved in''
  and credit the originating analysis.}
\end{frame}

\begin{frame}[shrink=7]{Why anyone should care}
  \begin{columns}[T]
    \begin{column}{0.52\textwidth}
      \textbf{1. Reproducible evaluation.}
      A score measured on a server under load can move when you repeat it on a
      quiet afternoon. Is that a regression, or did the numerical execution
      change?
      \vspace{0.5em}

      \textbf{2. Debugging.}
      You cannot bisect a failure you cannot reproduce. A bad generation that
      only happens under load is nearly impossible to chase.
      \vspace{0.5em}

      \textbf{3. Reinforcement learning}, which is where the problem stops
      being an annoyance and starts breaking the algorithm.
    \end{column}
    \begin{column}{0.48\textwidth}
      \begin{block}{Back to week 2}
        \small
        In GRPO you generate rollouts with an \key{inference engine} and
        compute the policy gradient with a \key{training engine}. The update
        assumes the log-probs match.
        \vspace{0.3em}
        They do not, partly for the reasons on the last three slides, so the
        RL you believed was on-policy is quietly off-policy.
        \vspace{0.3em}
        L3's fix was \key{truncated importance sampling}, multiplying the loss
        by a clipped ratio to correct the mismatch \cit{Yao et al. 2025}.
        \vspace{0.3em}
        \key{Batch-invariant kernels attack the same problem at the source.}
        If the engines agree bitwise you do not need the correction term, and
        the rollouts really are on-policy.
      \end{block}
    \end{column}
  \end{columns}
  \source{Yao et al. 2025 (truncated importance sampling), as covered in
    Lecture 3; adopted in Olmo 3's GRPO recipe.}
  \note{Land the three reasons in ascending order of seriousness and spend most
  of the time on the third.

  The RL connection is the payoff for the whole course structure, so make the
  link explicit. In week 2 I presented truncated importance sampling as one
  trick on a list of GRPO engineering fixes, justified in a line by
  ``inference-engine and training-engine log-probs differ''. Today they know
  why they differ.

  Say the deeper point slowly: a numerical detail in a matmul kernel changes
  the mathematical character of an RL algorithm. You cannot treat systems and
  algorithms as separate layers here, and that is a good reason for an NLP
  master's student to care about kernels.

  Note the DeepSeek V3.2 parallel if there is time. They log the MoE routing
  during rollout and force the same routing during training, which is the same
  instinct applied to a different mismatch.}
\end{frame}

\begin{frame}{Your turn: the Lab 3 bonus}
  \begin{columns}[T]
    \begin{column}{0.54\textwidth}
      \textbf{The experiment.}
      \begin{enumerate}\small
        \item Pick 50 prompts from your task.
        \item Send them at \texttt{temperature=0}, one at a time, and store the
          outputs.
        \item Send the same 50 again as a single large concurrent batch.
        \item \key{Count how many differ.} Find the first differing token
          position in each pair.
      \end{enumerate}
      \vspace{0.5em}
      \small
      Then the interesting follow-up: does the disagreement rate change with
      the \emph{length} of the output? With the batch size? With prefix
      caching on or off?
    \end{column}
    \begin{column}{0.46\textwidth}
      \begin{block}{What to write in the report}
        \small
        The count, the first-divergence positions, and one paragraph on what
        this implies for the evaluation numbers you report in weeks 3 and 4.
      \end{block}
      \vspace{0.3em}
      \begin{block}{Prediction to test}
        \small
        Short, confident outputs will mostly agree, and long generations and
        reasoning traces will not. If your task produces single-label answers
        you may see zero differences, which is a result too and tells you
        something about your margin.
      \end{block}
    \end{column}
  \end{columns}
  \vfill
  \centering
  \small
  You will believe this once you have measured it yourself, and not before.
  \note{Clock check: should be around 17:42.

  Sell this as the most genuinely fun thing in the lab. It takes twenty lines
  of code, it produces a surprising result, and it is a real open-ish problem
  rather than a reproduction exercise.

  Manage expectations honestly: students doing single-label classification may
  well see zero differences, because the margin between the top two logits is
  large. That is a good outcome and they should say so --- and then try it with
  a longer free-text prompt to see the effect.

  Encourage the follow-up questions in the small print. Anyone who plots
  disagreement rate against output length has done a small piece of genuine
  empirical work.}
\end{frame}
\fi

% =====================================================================
\iffalse
\begin{frame}[standout]
  Three things to remember
  \note{Pause here. Everything after this slide is summary, so this is the
  moment to stop adding information and start consolidating. Ask them to close
  laptops for ninety seconds. If the room is still with you, ask one person to
  name something they will change about how they run a model. The answers tell
  you which section landed.}
\end{frame}
\fi

\begin{frame}{Three things to remember}
  \begin{enumerate}
    \item \key{The two halves of a request hit different limits.} Prefill
      brings more arithmetic than the card can do at once, while decode
      finishes its arithmetic long before the weights have finished arriving.
      Find out which half you are paying for before you optimise, and
      remember that INT4 helps one of them and can hurt the other.
      \vspace{0.6em}
    \item \key{The KV cache often sets the concurrency limit.}
      In our worked example, eight gigabytes hold the weights and most of the
      remaining card is available to the cache. Paging, prefix reuse and
      quantisation all reduce pressure on that budget.
      \vspace{0.6em}
    \item \key{Measure latency and quality together.}
      Warm up the server, separate TTFT from TPOT, sweep concurrency, and
      report p50 and p95. A faster configuration is only useful if its output
      quality remains acceptable for the task.
  \end{enumerate}
  \vfill
  \textbf{Thursday:} you serve a model, measure it, quantise it, and
  find out how close today's estimates were.
  \note{These three sentences are reused verbatim in the week 7 wrap-up, so say
  them cleanly and do not improvise.

  If you are short of time, the one that must be said is number two, because it
  is the one that changes how they size a deployment. Number one is the
  explanation and number three is the memorable one.

  Then take thirty seconds on the hand-off slide. The restart script matters
  more to weeks 5 and 6 than anything else they will do on Thursday.}
\end{frame}

\begin{frame}{A server configuration you can reuse}
  \begin{columns}[T]
    \begin{column}{0.54\textwidth}
      If the same endpoint suits your project, it can become
      \key{infrastructure for later labs}:
      \vspace{0.4em}
      \begin{itemize}\small
        \item \textbf{Week 5, Dmitry's RAG lab}: your vLLM server is the
          \emph{generator} that retrieval feeds.
        \item \textbf{Week 6, the agent lab}: the same endpoint, now driving
          a tool-calling loop with many short calls, so TTFT and prefix caching
          suddenly matter a lot.
        \item \textbf{Week 7}: your demo runs against it.
      \end{itemize}
    \end{column}
    \begin{column}{0.46\textwidth}
      \begin{block}{The file worth keeping}
        \small
        \key{A one-command restart script in your repository.}
        \vspace{0.3em}
        Model path, LoRA adapter, \texttt{max-model-len},
        \texttt{max-num-seqs}, prefix caching and port, all of it in a file,
        so that in two weeks you type one line and the server you tuned today
        comes back exactly as it was.
      \end{block}
      \vspace{0.2em}
      \small
      The complete measurement workflow is in
      \texttt{lab3-example/README.md}.

      \warn{Every year someone loses an hour of week 5 rebuilding a server
      command from memory.} Do not be that person.
    \end{column}
  \end{columns}
  \vfill
  \centering
  \small
  GPU etiquette, again: the endpoint holds its KV pool whether or not you are
  using it. Shut it down when you leave.
  \note{Be firm about the script. It is one file and it takes two minutes,
  and without it week 5 starts with debugging instead of retrieval.

  Point out the week 6 detail, because an agent is a genuinely different
  serving regime: many short calls with a long shared prefix, which is the
  exact workload prefix caching was designed for. Their server configuration
  from Thursday may not be the right one for week 6, and noticing that is a
  good report paragraph.

  Final line, final thirty seconds: shut the server down. Then stop talking.}
\end{frame}

\appendix

\begin{frame}{Reading}
  \begin{columns}[T]
    \begin{column}{0.5\textwidth}
      \textbf{Serving systems}
      \begin{itemize}\footnotesize
        \item Kwon et al. 2023, \textit{PagedAttention} (vLLM)
        \item Yu et al. 2022, \textit{Orca}, continuous batching
        \item Agrawal et al. 2024, \textit{Sarathi-Serve}
        \item Zhong et al. 2024, \textit{DistServe}
        \item Zheng et al. 2024, \textit{SGLang} / RadixAttention
        \item Pope et al. 2023, \textit{Efficiently Scaling Transformer
          Inference}
      \end{itemize}
      \vspace{0.3em}
      \textbf{Kernels}
      \begin{itemize}\footnotesize
        \item Dao et al. 2022, Dao 2023, \textit{FlashAttention} 1 and 2
        \item Williams et al. 2009, the roofline model
      \end{itemize}
    \end{column}
    \begin{column}{0.5\textwidth}
      \textbf{Decoding}
      \begin{itemize}\footnotesize
        \item Leviathan et al. 2023, speculative decoding
        \item Chen et al. 2023, speculative sampling
        \item Cai et al. 2024, \textit{Medusa}; Li et al. 2024, \textit{EAGLE}
        \item Willard \& Louf 2023, \textit{Outlines}; Dong et al. 2024,
          \textit{XGrammar}
      \end{itemize}
      \vspace{0.3em}
      \textbf{Quantisation}
      \begin{itemize}\footnotesize
        \item Frantar et al. 2023, \textit{GPTQ}; Lin et al. 2024, \textit{AWQ}
        \item Hooper et al. 2024, \textit{KVQuant}
      \end{itemize}
      \vspace{0.3em}
      \textbf{Determinism}
      \begin{itemize}\footnotesize
        \item Thinking Machines Lab 2025, \textit{Defeating Nondeterminism
          in LLM Inference}
        \item SGLang deterministic-inference docs and issues
      \end{itemize}
    \end{column}
  \end{columns}
  \note{Venues, since the slide no longer has room for them: PagedAttention is
  SOSP 2023, Orca is OSDI 2022, Sarathi-Serve and DistServe are OSDI 2024,
  speculative decoding is ICML 2023, and Chen et al. is the DeepMind
  speculative-sampling report. The roofline paper is CACM 2009.

  If they only read three, say these. The vLLM paper, because it is the
  clearest systems paper in the area and the ideas come straight out of
  operating systems. The speculative decoding paper, for the proof, which is
  short and elegant. And the nondeterminism post, because it shows how to
  diagnose a problem everybody had learned to live with.

  For anyone who wants to work on this, the SGLang issue tracker is where the
  determinism discussion actually happens, and it is a realistic entry point
  for a thesis project. Say that I am happy to talk.}
\end{frame}

\end{document}
